\documentclass[12pt]{article}
\usepackage[a4paper,margin=1in]{geometry}
\usepackage{amsmath,amssymb,amsthm}
\usepackage{mathtools}
\usepackage{enumitem}
\usepackage{array}
\usepackage{microtype}
\usepackage[most]{tcolorbox}
\usepackage{xcolor}
\usepackage{setspace}
\usepackage{relsize}
\onehalfspacing
\usepackage{xcolor}
\usepackage[
colorlinks=true,
linkcolor=blue,
urlcolor=blue,
citecolor=blue
]{hyperref}


\DeclareMathOperator{\tr}{tr}
\DeclareMathOperator{\rank}{rank}

% ---------- Make big operators (sigma, product, etc.) render large ----------
\everymath{\displaystyle}
\let\oldsum\sum
\let\oldprod\prod
\let\oldbigcup\bigcup
\let\oldbigcap\bigcap
\let\oldint\int
\let\oldoplus\bigoplus
\renewcommand{\sum}{\mathlarger{\mathlarger{\oldsum}}}
\renewcommand{\prod}{\mathlarger{\mathlarger{\oldprod}}}
\renewcommand{\bigcup}{\mathlarger{\mathlarger{\oldbigcup}}}
\renewcommand{\bigcap}{\mathlarger{\mathlarger{\oldbigcap}}}
\renewcommand{\int}{\mathlarger{\mathlarger{\oldint}}}
\renewcommand{\bigoplus}{\mathlarger{\mathlarger{\oldoplus}}}

\definecolor{LightGreen}{RGB}{144,238,144}

% ---------- Every problem sits in its own light-green box ----------
\newtcolorbox[auto counter]{problem}{
	enhanced,
	breakable,
	colback=LightGreen!18,
	colframe=LightGreen!55!black,
	boxrule=0.8pt,
	arc=2pt,
	left=8pt,
	right=8pt,
	top=6pt,
	bottom=6pt,
	before skip=12pt,
	after skip=16pt,
	boxsep=2pt,
	before upper={\textbf{\thetcbcounter.}\quad}
}

\sloppy
\allowdisplaybreaks
\setlength{\jot}{10pt}
\clubpenalty=10000
\widowpenalty=10000

\title{\Huge\bfseries{Collection of Linear Algebra Problems}}
\author{}
\date{}

\begin{document}
\maketitle


\thispagestyle{empty}
\newpage
\begin{center}
	\vspace*{\fill}
	
	\textbf{Thanks to every Cat I have met in my life}
	
	
	\vspace*{\fill}
\end{center}


\newpage

\section*{A Note to the Reader}

This collection consists of 500 problems in linear algebra, carefully collected from a wide range of mathematical sources over the years. I collected problems from international and national competitions, mathematical journals and magazines, university-level contests, problem books, and other training materials.

Among the sources represented are the\textbf{International Mathematics Competition (IMC)}, \textbf{William Lowell Putnam Competetition}, \textbf{American Mathematical Monthly}, \textbf{Crux Mathematicorum}, \textbf{Romanian Mathematical Magazine},\textbf{Gazeta Mathematica}, \textbf{Romanian District Olympiad}, \textbf{Romanian National Mathematical Olympiad}, \textbf{Jozsef Wildt International Mathematical Competition}, \textbf{Vojtech Jarnik International Mathematical Competition}, \textbf{OLYMPIC Toán Sinh viên toàn quoc Lan}, \textbf{Berkeley Mathematics Competition}, and numerous other undergraduate mathematics competitions and problem collections.

I am deeply grateful to all the authors, problem setters, competition organisers, editors, members of AOPS, and other mathematical communities whose work made this collection possible. Their contributions have provided an invaluable source of challenging and interesting problems for students and enthusiasts of mathematics.

Since several problems have been translated from Romanian and Vietnamese with the assistance of Claude AI, there may be inaccuracies or ambiguities arising from translation or interpretation. Every effort has been made to preserve the mathematical meaning of the original problems; however, errors in translation, notation, or editing may still remain.
\medskip

If you notice any errors, inconsistencies, or omissions, or if you have a
correction or new problem suggestion for a future edition, please feel free to contact me at:

\noindent
\textbf{\texttt{category.theory.for.cats@gmail.com}}

\medskip

\noindent
\textbf{Online Repository}

\medskip

The source files and the continuously updated version of this collection are available on GitHub:

\medskip

\noindent
\href{https://github.com/ZetaZero-0/Collection-of-Problems}
{\textbf{Collection of Linear Algebra Problems}}

\vfill

\begin{center}
	\small\textit{Collection of Linear Algebra Problems}\\
	\small Author: \textit{Zeta Zero}\\
	\small Version: \textit{1.00}
\end{center}

\newpage

\begin{problem}
(Vandermonde determinant) Let
\[
V=\begin{pmatrix}
1 & x_1 & x_1^2 & \cdots & x_1^{n-1}\\
1 & x_2 & x_2^2 & \cdots & x_2^{n-1}\\
\vdots & \vdots & \vdots & \ddots & \vdots\\
1 & x_n & x_n^2 & \cdots & x_n^{n-1}
\end{pmatrix}.
\]
Show that $\det V = \displaystyle\prod_{1\le i<j\le n}(x_j-x_i)$.
\end{problem}

\begin{problem}
Let $A$ and $B$ be $n\times n$ matrices satisfying $A+B=AB$. Show that $AB=BA$.
\end{problem}

\begin{problem}
Suppose $A,B,C,D$ are $n\times n$ matrices, satisfying the conditions that $AB^T$ and $CD^T$ are symmetric and $AD^T-BC^T=I$. Prove that $A^TD-C^TB=I$.
\end{problem}

\begin{problem}
Do there exist square matrices $A$ and $B$ such that $AB-BA=I$?
\end{problem}

\begin{problem}
Let $A=(a_{ij})$ be a real $n\times n$ matrix satisfying $|a_{ii}| > \sum_{j\ne i}|a_{ij}|$ for all $1\le i\le n$. Prove that $A$ is invertible.
\end{problem}

\begin{problem}
Let
\[
A=\begin{pmatrix} 0 & 2 & -1\\ -2 & 0 & 2\\ 1 & -2 & 0\end{pmatrix}
\]
and let $I_3$ denote the $3\times 3$ identity matrix.
\begin{enumerate}[label=(\alph*)]
\item Verify that $A+I_3$ is nonsingular.
\item Verify that $(I_3-A)(I_3+A)^{-1} = (I_3+A)^{-1}(I_3-A)$.
\item Verify that the matrix $Q=(I_3-A)(I_3+A)^{-1}$ is orthogonal (i.e., $Q^T=Q^{-1}$).
\item Determine all real eigenvalues of the matrix $Q$, along with their corresponding eigenvectors.
\item Prove properties (a), (b), and (c) for the general case where $A$ is an arbitrary $3\times 3$ skew-symmetric matrix ($A^T=-A$).
\end{enumerate}
\end{problem}

\begin{problem}
Consider the quadratic form $q(x)=x^TAx$, where
\[
A=\begin{pmatrix} 2 & 1 & 0\\ 1 & 2 & 1\\ 0 & 1 & 2\end{pmatrix},
\]
and $x\in\mathbb{R}^3$ is a non-zero column vector.
\begin{enumerate}[label=(\alph*)]
\item Determine the orthogonal transformation $x=Qy$ that reduces the quadratic form to its canonical form.
\item Find the extreme values of the function $g:\mathbb{R}^3\setminus\{(0,0,0)^T\}\to\mathbb{R}$ given by $g(x)=\dfrac{x^TAx}{x^Tx}$, as well as the vectors where these values are attained.
\end{enumerate}
\end{problem}

\begin{problem}
Let $A\in M_2(\mathbb{C})$. Prove that for any positive integer $p\in\mathbb{N}^*$,
\[
\det(A^p+I_2) = (\det A)^p + \operatorname{tr}(A^p) + 1.
\]
\end{problem}

\begin{problem}
Consider a matrix $A\in M_2(\mathbb{C})$ such that $\det A=1$. Prove that
\[
\det(A^2+I_2) + \det(A^2+2A-I_2) = 8.
\]
\end{problem}

\begin{problem}
Let $A\in M_2(\mathbb{R})$ such that $\det A = d\ne 0$ and $\det(A+dA^*)=0$, where $A^*$ denotes the adjugate matrix of $A$. Prove that $\det(A-dA^*)=4$.
\end{problem}

\begin{problem}
Let $x>0$ be a real number and let $A\in M_2(\mathbb{R})$ be a matrix such that $\det(A^2+xI_2)=0$. Prove that $\det(A^2+A+xI_2)=x$.
\end{problem}

\begin{problem}
Let $A\in M_2(\mathbb{R})$. Prove that
\[
\det(A^2+A+I_2) \ge \frac{3}{4}(1-\det A)^2.
\]
\end{problem}

\begin{problem}
Let $a,b\in\mathbb{R}$ such that $b-a^2>0$. Find all matrices $A\in M_2(\mathbb{R})$ satisfying $\det(A^2-2aA+bI_2)=0$.
\end{problem}

\begin{problem}
Let $A\in M_3(\mathbb{C})$ be a matrix such that all entries on the main diagonals of $A$ and $A^2$ are zero. Prove that if $|\det A|<1$, then $\det(I_3+A^n)\ne 0$ for all $n\in\mathbb{N}^*$.
\end{problem}

\begin{problem}
Let $A\in M_3(\mathbb{C})$ be an invertible matrix. Prove that if $A^TA^{-1}=B$, then
\[
\det(I_3+B) = 2(1+\operatorname{tr}B) \quad\text{and}\quad \det(I_3+B^2) = 2(1+\operatorname{tr}B^2).
\]
\end{problem}

\begin{problem}
Let $A,B\in M_3(\mathbb{C})$ such that $(AB-BA)(AB-BA-I_3)=O_3$. Prove that $AB=BA$.
\end{problem}

\begin{problem}
Let $A\in M_4(\mathbb{R})$ be a symmetric matrix. Prove that $\det(A+I_4) = 4\sqrt{\det(A^2+I_4)}$ if and only if $A=I_4$.
\end{problem}

\begin{problem}
Let $A\in M_n(\mathbb{R})$. If $AA^T=I_n$, prove that $|\operatorname{tr}A|\le n$.
\end{problem}

\begin{problem}
Let $n\ge 2$ be an integer, and let $A,B,C\in M_n(\mathbb{C})$ be three mutually commuting matrices such that $A^{1999}=B^{1999}=C^{1999}=O_n$. Prove that the matrix $I_n+A^2+B^2+C^2$ is invertible.
\end{problem}

\begin{problem}
Let $A\in M_n(\mathbb{Q})$ be a matrix such that $\det\!\left(A+\sqrt[n]{2}\,I_n\right)=0$. Prove that
\[
\det(A-I_n) = \det A + \big(\det(A+I_n)\big)^n.
\]
\end{problem}

\begin{problem}
Let $A,B\in M_n(\mathbb{C})$ be two matrices such that $A^2-B^2 = 2(AB-BA)$. Prove that $(AB-BA)^n = O_n$.
\end{problem}

\begin{problem}
Let $A,B\in M_n(\mathbb{R})$ such that $B^2=I_n$ and $A^2=AB+I_n$. Prove that
\[
\det A \le \left(\frac{1+\sqrt5}{2}\right)^n.
\]
\end{problem}

\begin{problem}
Let $A\in M_n(\mathbb{C})$ be a matrix of rank $1$. Prove that for any complex number $\lambda$,
\[
\det(\lambda I_n+A) = \lambda^n + \lambda^{n-1}\operatorname{tr}A.
\]
\end{problem}

\begin{problem}
Let $A,X\in M_n(\mathbb{R})$ be matrices such that $\operatorname{rank}X=1$. Prove that
\[
\det(A+X)\det(A-X) \le (\det A)^2.
\]
\end{problem}

\begin{problem}
Let $n\ge 3$ be an odd integer, and let $A,B\in M_n(\mathbb{Z})$ be matrices such that $\det A\ne 0$ and $\operatorname{rank}B=1$. Prove that the following statements are equivalent:
\begin{enumerate}[label=(\roman*)]
\item $\det A = \operatorname{tr}B \in \{-1,1\}$;
\item $\det(A^*+BA^{-1}) - \det(A^*-A^{-1}B) = 2$, where $A^*$ is the adjugate matrix of $A$.
\end{enumerate}
\end{problem}

\begin{problem}
Let $\{a_k\}_{k=1}^\infty$ be a sequence of complex numbers, and let $b$ and $c$ also be complex numbers. For every positive integer $k$, define the $k\times k$ matrix $A_k$ by
\[
A_k := \begin{pmatrix}
a_1 & b & b & \cdots & b\\
c & a_2 & b & \cdots & b\\
c & c & a_3 & \cdots & b\\
\vdots & \vdots & \vdots & \ddots & \vdots\\
c & c & c & \cdots & a_k
\end{pmatrix}.
\]
\begin{enumerate}[label=\arabic*.]
\item When all $a_j=a$ for some complex number $a$ and when $b=c$, establish the following equation without using the formula in part 2 or the recursion in part 3: for $k\ge 1$, $\det(A_k) = (a-b)^{k-1}(a+kb-b)$.
\item When $k\ge 2$ and $b=c$, show that $\det(A_k) = 1 - b\dfrac{d}{db}\displaystyle\prod_{j=1}^k (a_j-b)$.
\item Show that for $k\ge 2$, $\det(A_{k+1}) = (a_k+a_{k+1}-b-c)\det(A_k) - (b-a_k)(c-a_k)\det(A_{k-1})$.
\end{enumerate}
\end{problem}

\begin{problem}
Let $A,B\in M_n(\mathbb{C})$ and $P\in\mathbb{C}[X]$ such that $A^n=O_n$ and $P(0)=0$. If $AB=BP(A)$, prove that $\det(A+B)=\det(B)$.
\end{problem}

\begin{problem}
Let $T_n=(t_{i,j})$ be the $n\times n$ matrix with $t_{i,j}=\tan\big((i+j-1)x\big)$, i.e.,
\[
T_n = \begin{pmatrix}
\tan x & \tan 2x & \cdots & \tan nx\\
\tan 2x & \tan 3x & \cdots & \tan(n+1)x\\
\vdots & \vdots & \ddots & \vdots\\
\tan nx & \tan(n+1)x & \cdots & \tan(2n-1)x
\end{pmatrix}.
\]
Prove that $\det(T_n) = (-1)^{\lfloor n/2\rfloor}\sec^n(nx)$.
\end{problem}

\begin{problem}
Let $n$ be a positive integer and let $A$ be the $n\times n$ matrix defined by $A_{i,j}=1$ if $\min(i,j)$ is odd, $2$ if $\min(i,j)$ is even, i.e.,
\[
A = \begin{pmatrix}
1 & 1 & 1 & 1 & \cdots & 1\\
1 & 2 & 2 & 2 & \cdots & 2\\
1 & 2 & 1 & 1 & \cdots & 1\\
1 & 2 & 1 & 2 & \cdots & 2\\
\vdots & \vdots & \vdots & \vdots & \ddots & \vdots
\end{pmatrix}.
\]
Show that $A^{-1}$ exists and determine the number of zero entries in $A^{-1}$.
\end{problem}

\begin{problem}
Given two complex $n\times n$ positive-definite matrices $A$ and $B$, let
\[
C=\tfrac12(A+B), \qquad D = A^{1/2}\big(A^{-1/2}BA^{-1/2}\big)^{1/2}A^{1/2};
\]
the matrices $C$ and $D$ are the arithmetic mean and the geometric mean of $A$ and $B$. Prove that
\begin{enumerate}[label=(\alph*)]
\item $\operatorname{range}(C-D) = \operatorname{range}(A-B)$.
\item $\operatorname{range}\begin{pmatrix}C & D\\ D & C\end{pmatrix} = \operatorname{range}\begin{pmatrix}A & B\\ B & A\end{pmatrix}$.
\end{enumerate}
\end{problem}

\begin{problem}
Let $a_i\in\mathbb{R}$ and let $b_i\in\mathbb{R}$. Define the $n\times n$ matrix
\[
M = \begin{pmatrix}
a_1b_1 & a_1b_2 & a_1b_3 & \cdots & a_1b_n\\
a_2b_1 & a_2b_2 & a_2b_3 & \cdots & a_2b_n\\
a_3b_1 & a_3b_2 & a_3b_3 & \cdots & a_3b_n\\
\vdots & \vdots & \vdots & \ddots & \vdots\\
a_nb_1 & a_nb_2 & a_nb_3 & \cdots & a_nb_n
\end{pmatrix}.
\]
What is the maximum rank of $M$? When does it attain maximum rank?
\end{problem}

\begin{problem}
An $n\times n$ matrix $A$ with real entries is called interesting if
\[
A^2 = 2A^T - I \quad\text{and}\quad \operatorname{tr}(A)+\operatorname{tr}(A^2) = 0.
\]
\begin{enumerate}[label=\alph*.]
\item Find all positive integers $n$ such that an interesting matrix of size $n\times n$ exists.
\item Find the possible determinants of $A$ for each positive integer $n$ for which interesting matrices of size $n\times n$ exist.
\end{enumerate}
\end{problem}

\begin{problem}
Let $n\ge 3$. Let $A$ be the $n\times n$ matrix with $A_{jk}=\cos(2\pi(j+k)/n)$. Find $\det(I+A)$.
\end{problem}

\begin{problem}
Let $A$ and $B$ be $n\times n$ matrices. Show that $\det(I+AB) = \det(I+BA)$.
\end{problem}

\begin{problem}
Let $A,B,C,D$ be $n\times n$ matrices such that $AC=CA$. Prove that
\[
\det\begin{pmatrix}A & B\\ C & D\end{pmatrix} = \det(AD-CB).
\]
\end{problem}

\begin{problem}
Let $A$ be a square matrix with real entries. Show that $\det(A^2+I)\ge 0$.
\end{problem}

\begin{problem}
Let $A,B,C$ be $n\times n$ real matrices that pairwise commute and $ABC=O$. Show that $\det(A^3+B^3+C^3)\det(A+B+C) \ge 0$.
\end{problem}

\begin{problem}
Let $A$ and $B$ be two $n\times n$ real matrices that commute. Suppose that $\det(A+B)\ge 0$. Prove that $\det(A^k+B^k) \ge 0$ for all $k\ge 1$.
\end{problem}

\begin{problem}
Let $A$ be a real skew-symmetric square matrix (i.e., $A^T=-A$). Prove that $\det(I+tA^2) \ge 0$ for all real $t$.
\end{problem}

\begin{problem}
Let $A$ and $B$ be real symmetric matrices with all eigenvalues strictly greater than $1$. Let $\lambda$ be a real eigenvalue of matrix $AB$. Prove that $|\lambda|>1$.
\end{problem}

\begin{problem}
Let $V$ be a $10$-dimensional real vector space and $U_1$ and $U_2$ two linear subspaces such that $U_1\subseteq U_2$, $\dim_{\mathbb{R}} U_1=3$ and $\dim_{\mathbb{R}} U_2=6$. Let $E$ be the set of all linear maps $T:V\to V$ which have $U_1$ and $U_2$ as invariant subspaces (i.e. $T(U_1)\subseteq U_1$ and $T(U_2)\subseteq U_2$). Calculate the dimension of $E$ as a real vector space.
\end{problem}

\begin{problem}
Determine all pairs $(a,b)$ of real numbers for which there exists a unique symmetric $2\times 2$ matrix $M$ with real entries satisfying $\operatorname{trace}(M)=a$ and $\det(M)=b$.
\end{problem}

\begin{problem}
For any integer $n>2$ and two $n\times n$ matrices with real entries $A,B$ that satisfy the equation $A^{-1}+B^{-1}=(A+B)^{-1}$, prove that $\det(A)=\det(B)$. Does the same conclusion follow for matrices with complex entries?
\end{problem}

\begin{problem}
\begin{enumerate}[label=(\alph*)]
\item Show that for any $m\in\mathbb{N}$ there exists a real $m\times m$ matrix $A$ such that $A^3=A+I$ where $I$ is the $m\times m$ identity matrix.
\item Show that $\det(A)>0$ for every real $m\times m$ matrix satisfying $A^3=A+I$.
\end{enumerate}
\end{problem}

\begin{problem}
Let $A$ be the $n\times n$ matrix, whose $(i,j)$-th entry is $i+j$ for all $i,j=1,2,\dots,n$. What is the rank of $A$?
\end{problem}

\begin{problem}
Determine all complex numbers $\lambda$ for which there exist a positive integer $n$ and a real $n\times n$ matrix $A$ such that $A^2=A^T$ and $\lambda$ is an eigenvalue of $A$.
\end{problem}

\begin{problem}
Let $A$ be a real $n\times n$ matrix such that $A^3=0$.
\begin{enumerate}[label=(\alph*)]
\item Prove that there is a unique real $n\times n$ matrix $X$ that satisfies the equation $X+AX+XA^2=A$.
\item Express $X$ in terms of $A$.
\end{enumerate}
\end{problem}

\begin{problem}
\begin{enumerate}[label=(\alph*)]
\item Let $A$ be an $n\times n$, $n\ge 2$, symmetric invertible matrix with real positive elements. Show that $z_n \le n^2-2n$, where $z_n$ is the number of zero elements in $A^{-1}$.
\item How many zero elements are there in the inverse of the $n\times n$ matrix
\[
A = \begin{pmatrix}
1 & 1 & 1 & \cdots & 1\\
1 & 2 & 2 & \cdots & 2\\
1 & 2 & 1 & \cdots & 1\\
\vdots & \vdots & \vdots & \ddots & \vdots
\end{pmatrix}\ ?
\]
\end{enumerate}
\end{problem}

\begin{problem}
Let $V$ be a real vector space, and let $f,f_1,f_2,\dots,f_k$ be linear maps from $V$ to $\mathbb{R}$. Suppose that $f(x)=0$ whenever $f_1(x)=f_2(x)=\cdots=f_k(x)=0$. Prove that $f$ is a linear combination of $f_1,f_2,\dots,f_k$.
\end{problem}

\begin{problem}
Let $A$ be a $3\times 3$ real matrix such that the vectors $Au$ and $u$ are orthogonal for each column vector $u\in\mathbb{R}^3$. Prove that:
\begin{enumerate}[label=(\alph*)]
\item $A^T=-A$, where $A^T$ denotes the transpose of $A$;
\item there exists a vector $v\in\mathbb{R}^3$ such that $Au = v\times u$ for every $u\in\mathbb{R}^3$, where $v\times u$ denotes the vector product in $\mathbb{R}^3$.
\end{enumerate}
\end{problem}

\begin{problem}
Let $a_j=a_0+jd$ for $j=0,\dots,n$ where $a_0,d$ are fixed real numbers. Put
\[
A = \begin{pmatrix}
a_0 & a_1 & a_2 & \cdots & a_n\\
a_1 & a_0 & a_1 & \cdots & a_{n-1}\\
\vdots & \vdots & \vdots & \ddots & \vdots\\
a_n & a_{n-1} & a_{n-2} & \cdots & a_0
\end{pmatrix}.
\]
Calculate $\det(A)$.
\end{problem}

\begin{problem}
Compute the determinant of the $n\times n$ matrix $A=[a_{ij}]$:
\[
a_{ij} = \begin{cases} (-1)^{|i-j|}, & i\ne j\\ 2, & i=j.\end{cases}
\]
\end{problem}

\begin{problem}
Let $A$ be a real $4\times 2$ matrix and $B$ be a real $2\times 4$ matrix such that
\[
AB = \begin{pmatrix} 1 & 0 & -1 & 0\\ 0 & 1 & 0 & -1\\ -1 & 0 & 1 & 0\\ 0 & -1 & 0 & 1\end{pmatrix}.
\]
Find $BA$.
\end{problem}

\begin{problem}
Let $k$ be a positive integer. Find the smallest positive integer $n$ for which there exist $k$ nonzero vectors $v_1,\dots,v_k$ in $\mathbb{R}^n$ such that for every pair $i,j$ of indices with $|i-j|>1$ the vectors $v_i$ and $v_j$ are orthogonal.
\end{problem}

\begin{problem}
Let $A$ be a real $n\times n$ matrix and suppose that for every positive integer $m$ there exists a real symmetric matrix $B$ such that $2021B = A^m+B^2$. Prove that $|\det A| \le 1$.
\end{problem}

\begin{problem}
Let $n$ be a positive integer. Find all $n\times n$ real matrices $A$ with only real eigenvalues satisfying $A+A^k = A^T$ for some integer $k\ge n$.
\end{problem}

\begin{problem}
Let $n\ge 2$ be an integer. What is the minimal and maximal possible rank of an $n\times n$ matrix whose $n^2$ entries are precisely the numbers $1,2,\dots,n^2$?
\end{problem}

\begin{problem}
Let $M$ be an invertible matrix of dimension $2n\times 2n$, represented in block form as
\[
M = \begin{pmatrix} A & B\\ C & D\end{pmatrix}, \qquad M^{-1} = \begin{pmatrix} E & F\\ G & H\end{pmatrix}.
\]
Show that $\det M \cdot \det H = \det A$.
\end{problem}

\begin{problem}
Let $A$ and $B$ be $n\times n$ real matrices such that $\operatorname{rank}(AB-BA+I)=1$. Prove that
\[
\operatorname{trace}(ABAB) - \operatorname{trace}(A^2B^2) = \tfrac12 n(n-1).
\]
\end{problem}

\begin{problem}
A four-digit number YEAR is called very good if the system
\[
\begin{aligned}
Yx+Ey+Az+Rw &= Y\\
Rx+Yy+Ez+Aw &= E\\
Ax+Ry+Yz+Ew &= A\\
Ex+Ay+Rz+Yw &= R
\end{aligned}
\]
of linear equations in the variables $x,y,z,w$ has at least two solutions. Find all very good YEARs in the 21st century.
\end{problem}

\begin{problem}
Let $A=(a_{ij})_{i,j=1}^n$ be a symmetric $n\times n$ matrix with real entries, and let $\lambda_1,\lambda_2,\dots,\lambda_n$ denote its eigenvalues. Show that
\[
\sum_{1\le i<j\le n} a_{ii}a_{jj} \ge \sum_{1\le i<j\le n}\lambda_i\lambda_j
\]
and determine all matrices for which equality holds.
\end{problem}

\begin{problem}
Let $n$ be a fixed positive integer. Determine the smallest possible rank of an $n\times n$ matrix that has zeros along the main diagonal and strictly positive real numbers off the main diagonal.
\end{problem}

\begin{problem}
Does there exist a real $3\times 3$ matrix $A$ such that $\operatorname{tr}(A)=0$ and $A^2+A^T=I$?
\end{problem}

\begin{problem}
Let $A_1,A_2,\dots,A_k$ be $n\times n$ idempotent complex matrices such that $A_iA_j=-A_jA_i$ for all $i\ne j$. Prove that at least one of the given matrices has rank $\le n/k$.
\end{problem}

\begin{problem}
Determine all rational numbers $a$ for which the matrix
\[
\begin{pmatrix} a & -a & -1 & 0\\ a & -a & 0 & -1\\ 1 & 0 & a & -a\\ 0 & 1 & a & -a\end{pmatrix}
\]
is the square of a matrix with all rational entries.
\end{problem}

\begin{problem}
Define the sequence $A_1,A_2,\dots$ of matrices by the following recurrence:
\[
A_1 = \begin{pmatrix}0 & 1\\ 1 & 0\end{pmatrix}, \qquad A_{n+1} = \begin{pmatrix} A_n & I_{2^n}\\ I_{2^n} & A_n\end{pmatrix} \quad (n=1,2,\dots)
\]
where $I_m$ is the $m\times m$ identity matrix. Prove that $A_n$ has $n+1$ distinct integer eigenvalues $\lambda_0<\lambda_1<\cdots<\lambda_n$ with multiplicities $\binom{n}{0},\binom{n}{1},\dots,\binom{n}{n}$, respectively.
\end{problem}

\begin{problem}
Let $A,B\in M_n(\mathbb{C})$ be two $n\times n$ matrices such that
\[
A^2B+BA^2 = 2ABA.
\]
Prove that there exists a positive integer $k$ such that $(AB-BA)^k=0$.
\end{problem}

\begin{problem}
Let $n$ be a positive integer, and consider the matrix $A=(a_{ij})_{1\le i,j\le n}$, where
\[
a_{ij} = \begin{cases} 1 & \text{if } i+j \text{ is a prime number}\\ 0 & \text{otherwise.}\end{cases}
\]
Prove that $|\det A| = k^2$ for some integer $k$.
\end{problem}

\begin{problem}
Let $n>1$ be an odd positive integer and $A=(a_{ij})_{i,j=1\dots n}$ be the $n\times n$ matrix with
\[
a_{ij} = \begin{cases} 2 & \text{if } i=j\\ 1 & \text{if } i-j\equiv \pm 2 \pmod n\\ 0 & \text{otherwise.}\end{cases}
\]
Find $\det(A)$.
\end{problem}

\begin{problem}
Let $A$ be an $n\times n$-matrix with integer entries and $b_1,\dots,b_k$ be integers satisfying $\det A = b_1\cdots b_k$. Prove that there exist $n\times n$ matrices $B_1,\dots,B_k$ with integer entries such that $A=B_1\cdots B_k$ and $\det B_i = b_i$ for all $i=1,\dots,k$.
\end{problem}

\begin{problem}
Let $v_0$ be the zero vector in $\mathbb{R}^n$ and let $v_1,v_2,\dots,v_{n+1}\in\mathbb{R}^n$ be such that the Euclidean norm $|v_i-v_j|$ is rational for every $0\le i,j\le n+1$. Prove that $v_1,\dots,v_{n+1}$ are linearly dependent over the rationals.
\end{problem}

\begin{problem}
In the linear space of all real $n\times n$ matrices, find the maximum possible dimension of a linear subspace $V$ such that $\forall X,Y\in V$, $\operatorname{trace}(XY)=0$.
\end{problem}

\begin{problem}
For $n\ge 1$ let $M$ be an $n\times n$ complex matrix with distinct eigenvalues $\lambda_1,\lambda_2,\dots,\lambda_k$ with multiplicities $m_1,m_2,\dots,m_k$ respectively. Consider the linear operator $L_M$ defined by $L_M(X)=MX+XM^T$, for any complex $n\times n$ matrix $X$. Find its eigenvalues and their multiplicities.
\end{problem}

\begin{problem}
Let $A$ be an $n\times n$ real matrix such that $3A^3 = A^2+A+I$. Show that the sequence $A^k$ converges to an idempotent matrix.
\end{problem}

\begin{problem}
Let $A,B$ be square complex matrices of the same size such that $\operatorname{rank}(AB-BA)=1$. Show that $(AB-BA)^2=0$.
\end{problem}

\begin{problem}
Let $\alpha\in\mathbb{R}\setminus\{0\}$ and suppose that $F$ and $G$ are linear maps (operators) from $\mathbb{R}^n$ into $\mathbb{R}^n$ satisfying $F\circ G - G\circ F = \alpha F$.
\begin{enumerate}[label=(\alph*)]
\item Show that for all $k\in\mathbb{N}$ one has $F^k\circ G - G\circ F^k = \alpha k F^k$.
\item Show that there exists $k\ge 1$ such that $F^k=0$.
\end{enumerate}
\end{problem}

\begin{problem}
Let $A$ be $n\times n$ diagonal matrix with characteristic polynomial
\[
(x-c_1)^{d_1}(x-c_2)^{d_2}\cdots(x-c_k)^{d_k}
\]
where $c_1,c_2,\dots,c_k$ are distinct. Let $V$ be the space of all $n\times n$ matrices $B$ such that $AB=BA$. Prove that the dimension of $V$ is $d_1^2+d_2^2+\cdots+d_k^2$.
\end{problem}

\begin{problem}
Let $n$ be a positive integer. At most how many distinct unit vectors can be selected in $\mathbb{R}^n$ such that from any three of them, at least two are orthogonal?
\end{problem}

\begin{problem}
Let $A$ be a symmetric $m\times m$ matrix over the two-element field all of whose diagonal entries are zero. Prove that for every positive integer $n$ each column of the matrix $A^n$ has a zero entry.
\end{problem}

\begin{problem}
Let $A$ and $B$ be real $n\times n$ matrices. Assume that there exist $n+1$ different real numbers $t_1,t_2,\dots,t_{n+1}$ such that the matrices $C_i=A+t_iB$, $i=1,2,\dots,n+1$ are nilpotent (i.e. $C_i^n=0$). Show that both $A$ and $B$ are nilpotent.
\end{problem}

\begin{problem}
The linear operator $A$ on the vector space $V$ is called an involution if $A^2=E$ where $E$ is the identity operator on $V$. Let $\dim V = n<\infty$.
\begin{enumerate}[label=(\roman*)]
\item Prove that for every involution $A$ on $V$ there exists a basis of $V$ consisting of eigenvectors of $A$.
\item Find the maximal number of distinct pairwise commuting involutions on $V$.
\end{enumerate}
\end{problem}

\begin{problem}
Determine all positive integers $n$ for which there exist $n\times n$ real invertible matrices $A$ and $B$ that satisfy $AB-BA = B^2A$.
\end{problem}

\begin{problem}
\begin{enumerate}[label=(\alph*)]
\item Let $f:M_n\to\mathbb{R}$ be a linear map from the space $M_n=\mathbb{R}^{n^2}$ of real $n\times n$ matrices to the reals. Prove that there exists a unique matrix $C\in M_n$ such that $f(A)=\operatorname{tr}(AC)$ for any $A\in M_n$.
\item Suppose in addition to (a) that $f(AB)=f(BA)$ for any $A,B\in M_n$. Prove that there exists $\lambda\in\mathbb{R}$ such that $f(A)=\lambda\cdot\operatorname{tr}(A)$.
\end{enumerate}
\end{problem}

\begin{problem}
An $n\times n$ complex matrix $A$ is called t-normal if $AA^T=A^TA$ where $A^T$ is the transpose of $A$. For each $n$, determine the maximum dimension of a linear space of complex $n\times n$ matrices consisting of t-normal matrices.
\end{problem}

\begin{problem}
Let $A$ be an $n\times n$ complex matrix whose eigenvalues have absolute value at most $1$. Prove that $\|A^n\| \le n^{\ln 2}\|A\|^{n-1}$. (Here $\|B\| = \sup_{\|x\|\le 1}\|Bx\|$ for every $n\times n$ matrix $B$ and $\|x\| = \sqrt{\sum_{i=1}^n |x_i|^2}$ for every complex vector $x\in\mathbb{C}^n$.)
\end{problem}

\begin{problem}
Let $A$ be an $n\times n$ matrix with complex entries and suppose that $n>1$. Prove that
\[
A\overline{A} = I_n \iff \exists\, S\in GL_n(\mathbb{C}) \text{ such that } A = S^{-1}\overline{S}.
\]
\end{problem}

\begin{problem}
Determine whether there exist an odd positive integer $n$ and $n\times n$ matrices $A$ and $B$ with integer entries, that satisfy the following conditions:
\begin{enumerate}[label=(\alph*)]
\item $\det(B)=1$
\item $AB=BA$
\item $A^4+4A^2B^2+16B^4 = 2019I$
\end{enumerate}
\end{problem}

\begin{problem}
For each positive integer $k$, find the smallest number $n_k$ for which there exist real $n_k\times n_k$ matrices $A_1,A_2,\dots,A_k$ such that all of the following conditions hold:
\begin{enumerate}[label=(\arabic*)]
\item $A_1^2=A_2^2=\cdots=A_k^2=0$,
\item $A_iA_j=A_jA_i$ for all $1\le i,j\le k$, and
\item $A_1A_2\cdots A_k \ne 0$.
\end{enumerate}
\end{problem}

\begin{problem}
For an $m\times m$ real matrix $A$, $e^A$ is defined as $\sum_{n=0}^\infty \frac1{n!}A^n$. (The sum is convergent for all matrices.) Prove or disprove that for all real polynomials $p$ and $m\times m$ real matrices $A$ and $B$, $p(e^AB)$ is nilpotent if and only if $p(e^BA)$ is nilpotent.
\end{problem}

\begin{problem}
For $n\ge 0$ define matrices $A_n$ and $B_n$ as follows: $A_0=B_0=(1)$ and for every $n>0$,
\[
A_n = \begin{pmatrix} A_{n-1} & A_{n-1}\\ A_{n-1} & B_{n-1}\end{pmatrix}, \qquad B_n = \begin{pmatrix} A_{n-1} & A_{n-1}\\ A_{n-1} & 0\end{pmatrix}.
\]
Denote the sum of all elements of a matrix $M$ by $S(M)$. Prove that $S(A_n^{k-1}) = S(A_k^{n-1})$ for every $n,k\ge 1$.
\end{problem}

\begin{problem}
Let $A_i,B_i,S_i$ ($i=1,2,3$) be invertible real $2\times 2$ matrices such that:
\begin{enumerate}[label=(\arabic*)]
\item not all $A_i$ have a common real eigenvector;
\item $A_i = S_i^{-1}B_iS_i$ for all $i=1,2,3$;
\item $A_1A_2A_3 = B_1B_2B_3 = \begin{pmatrix}1 & 0\\ 0 & 1\end{pmatrix}$.
\end{enumerate}
Prove that there is an invertible real $2\times 2$ matrix $S$ such that $A_i=S^{-1}B_iS$ for all $i=1,2,3$.
\end{problem}

\begin{problem}
Let $A^{n\times n}$ and $B^{n\times n}$ be two matrices which differ only by one row/column. Show that $\det(A+B) = 2^{n-1}\{\det(A)+\det(B)\}$.
\end{problem}

\begin{problem}
Let $A=(a_{ij})$ and $B=(b_{ij})$ be two square matrices of the same dimension.
\begin{enumerate}[label=(\alph*)]
\item Let $a_{ij}=\rho^{|i-j|}$; show that $\det(A) = (1-\rho^2)^{n-1}$.
\item Suppose now $b_{ij}=\rho^{i-j}a_{ij}$. Find $\det(B)$ in terms of $\rho$.
\end{enumerate}
\end{problem}

\begin{problem}
Consider the matrix
\[
B^{n\times n}(x) = \begin{pmatrix}
1 & x & 0 & \cdots & 0\\
x & 1 & x & \cdots & 0\\
0 & x & 1 & \cdots & 0\\
\vdots & \vdots & \vdots & \ddots & \vdots\\
0 & 0 & 0 & \cdots & 1
\end{pmatrix}, \quad x\in\mathbb{R};
\]
find $\det(B^{n\times n}(x))$ for any $x\in\mathbb{R}$. What if $x=\tfrac12$?
\end{problem}

\begin{problem}
Prove that
\[
\begin{vmatrix}
1+a_1 & a_2 & \cdots & a_n\\
a_1 & 1+a_2 & \cdots & a_n\\
\vdots & \vdots & \ddots & \vdots\\
a_1 & a_2 & \cdots & 1+a_n
\end{vmatrix} = 1+\sum_{i=1}^n a_i.
\]
\end{problem}

\begin{problem}
Let $g(x)=\prod_{i=1}^n (u_i-x)$ and let
\[
\Delta_n = \begin{vmatrix}
u_1 & a & a & \cdots & a\\
b & u_2 & a & \cdots & a\\
b & b & u_3 & \cdots & a\\
\vdots & \vdots & \vdots & \ddots & a\\
b & b & b & \cdots & u_n
\end{vmatrix}.
\]
\begin{enumerate}[label=(\alph*)]
\item Show that if $a\ne b$ then $\Delta_n = \dfrac{bg(a)-ag(b)}{b-a}$.
\item Also show that if $a=b$ then $\Delta_n = a\sum_{i=1}^n g_i(a) + u_ng_n(a)$ where $g_i(a) = \dfrac{g(a)}{u_i-a}$ for $i=1,\dots,n$.
\item Use (b) to evaluate $\det\big((a-b)I_n + b\,\mathbf 1_n\mathbf 1_n^T\big)$.
\end{enumerate}
\end{problem}

\begin{problem}
Show that if $a\ne b$ then
\[
\begin{vmatrix}
a+b & ab & 0 & 0 & \cdots & 0\\
1 & a+b & ab & 0 & \cdots & 0\\
0 & 1 & a+b & ab & \cdots & 0\\
0 & 0 & 1 & a+b & \cdots & 0\\
\vdots & \vdots & \vdots & \vdots & \ddots & \vdots\\
0 & 0 & 0 & \cdots & 1 & a+b
\end{vmatrix} = \frac{a^{n+1}-b^{n+1}}{a-b};
\]
what if $a=b$?
\end{problem}

\begin{problem}
Let $p<n$ be any two positive integers. Prove that $\det(\Delta_{n,p})=1$ where $\Delta_{n,p}$ is the matrix with $(i,j)$-th entry $\delta_{ij}^{n,p} = \binom{n+i-1}{j-1}$ for $1\le i,j\le p+1$.
\end{problem}

\begin{problem}
Prove that the determinant of a circulant matrix is of the following form:
\[
\begin{vmatrix}
c_1 & c_2 & c_3 & \cdots & c_n\\
c_n & c_1 & c_2 & \cdots & c_{n-1}\\
\vdots & \vdots & \vdots & \ddots & \vdots\\
c_2 & c_3 & c_4 & \cdots & c_1
\end{vmatrix} = (-1)^{n-1}\prod_{k=1}^n \left(\sum_{j=1}^n \zeta^{k(j-1)}c_j\right)
\]
where $\zeta = e^{2\pi i/n}$.
\end{problem}

\begin{problem}
Given distinct integers $x_1,x_2,\dots,x_n$, prove that $\prod_{i>j}(x_i-x_j)$ is divisible by $\prod_{k=1}^{n-1} k!$.
\end{problem}

\begin{problem}
Let $f_{ij}(t)$ be a differentiable function of $t$. Then show that
\[
\frac{d}{dt}\begin{vmatrix} f_{11}(t) & \cdots & f_{1n}(t)\\ \vdots & \ddots & \vdots\\ f_{n1}(t) & \cdots & f_{nn}(t)\end{vmatrix} = \sum_{j=1}^n \det(F_j(t))
\]
where $F_j(t)$ is the matrix with the $j$-th row differentiated.
\end{problem}

\begin{problem}
If $A$ is an $n\times n$ matrix over $\mathbb{C}$ then show that $|\det(A)| \le \prod_{i=1}^n\left(\sum_{j=1}^n |a_{ij}|\right)$.
\end{problem}

\begin{problem}
Show that the determinant of the Cauchy matrix $A=(a_{ij})$ where $a_{ij}=\dfrac{1}{x_i+y_j}$ is of the following form:
\[
\det(A) = \frac{\prod_{i>j}(x_i-x_j)(y_i-y_j)}{\prod_{i,j}(x_i+y_j)}.
\]
\end{problem}

\begin{problem}
(The Binet--Cauchy Formula) Let $A$ and $B$ be matrices of size $n\times m$ and $m\times n$ where $n\le m$. Then show that
\[
\det(AB) = \sum_{1\le k_1<k_2<\cdots<k_n\le m} A^{k_1,k_2,\dots,k_n}B_{k_1,k_2,\dots,k_n}
\]
where $A^{k_1,\dots,k_n}$ is the minor obtained from $A$ using columns $k_1,\dots,k_n$ and $B_{k_1,\dots,k_n}$ is the minor obtained from $B$ using rows $k_1,\dots,k_n$.
\end{problem}

\begin{problem}
Let $M$ be an $n\times n$ matrix such that $M+M^T=O_n$. Prove that for all $\epsilon\in\mathbb{R}$, $\det(I_n+\epsilon M^2) \ge 0$.
\end{problem}

\begin{problem}
Let $S_i = \sum_{j=1}^n b_{ij}$. Prove that
\[
\begin{vmatrix} S_1-b_{11} & \cdots & S_1-b_{1n}\\ \vdots & \ddots & \vdots\\ S_n-b_{n1} & \cdots & S_n-b_{nn}\end{vmatrix} = (-1)^{n-1}(n-1)\det(B)
\]
where $B=(b_{ij})_{i,j=1}^n$.
\end{problem}

\begin{problem}
Prove that
\[
\begin{vmatrix}
a_{11} & \cdots & a_{1n} & b_1\\
\vdots & \ddots & \vdots & \vdots\\
a_{n1} & \cdots & a_{nn} & b_n\\
c_1 & \cdots & c_n & 0
\end{vmatrix} = -\sum_{i,j} b_ic_jA_{ij}
\]
where $A_{ij}$ is the cofactor of $a_{ij}$ in $A$.
\end{problem}

\begin{problem}
Let $u$ and $v$ be columns of length $n$. Prove that $\operatorname{adj}(I-uv^T) = uv^T + (1-v^Tu)I$.
\end{problem}

\begin{problem}
Let $A,B,C,D$ be $n\times n$ matrices such that $ABCD=I$. Show that $ABCD = DABC = CDAB = BCDA = I$.
\end{problem}

\begin{problem}
Find the inverse of the matrix $A=(a_{ij})_{i,j=1}^n$ where $a_{ij}=\min(i,j)$ for $i,j=1,\dots,n$.
\end{problem}

\begin{problem}
Let $A,B,C,D\in M_{n\times n}(\mathbb{R})$ and the matrix $\begin{pmatrix}A & B\\ C & D\end{pmatrix}$ has rank $n$. Show that
\[
\left|\begin{matrix}\det(A) & \det(B)\\ \det(C) & \det(D)\end{matrix}\right| = 0;
\]
also if $A$ is invertible, then $D=CA^{-1}B$.
\end{problem}

\begin{problem}
Let $A$ and $B$ be $n\times n$ complex matrices. Show that
\[
\left|\begin{matrix} A & B\\ B & A\end{matrix}\right| = \det(A+B)\det(A-B).
\]
Let $E=A+B$ and $F=A-B$. If $E$ and $F$ are invertible then
\[
\begin{pmatrix}A & B\\ B & A\end{pmatrix}^{-1} = \frac12\begin{pmatrix} E^{-1}+F^{-1} & E^{-1}-F^{-1}\\ E^{-1}-F^{-1} & E^{-1}+F^{-1}\end{pmatrix}.
\]
\end{problem}

\begin{problem}
Let $A\in M_n(\mathbb{R})$ be a skew-symmetric matrix. Show that $\det(A)\ge 0$.
\end{problem}

\begin{problem}
Show that if $A,B\in M_n(\mathbb{R})$ such that $AB=BA$ then $\det(A^2+B^2)\ge 0$.
\end{problem}

\begin{problem}
Let $A,B\in M_2(\mathbb{R})$ such that $AB=BA$ and $\det(A^2+B^2)=0$. Prove that $\det(A)=\det(B)$.
\end{problem}

\begin{problem}
Let $A\in M_2(\mathbb{R})$ such that $\det(A)=-1$. Show that $\det(A^2+I_2)\ge 4$. When does the equality hold?
\end{problem}

\begin{problem}
Let $A,B\in M_3(\mathbb{C})$ such that $\det(A)=\det(B)=1$. Show that $\det(A+\sqrt2\,B)\ne 0$.
\end{problem}

\begin{problem}
Let $A,B\in M_2(\mathbb{R})$. Show that $\det\big((AB+BA)^4+(AB-BA)^4\big) \ge 0$.
\end{problem}

\begin{problem}
\begin{enumerate}[label=(\alph*)]
\item Let $A$ be a matrix from $M_2(\mathbb{C})$, $A\ne aI_2$ for any $a\in\mathbb{C}$. Prove that the matrix $X$ from $M_2(\mathbb{C})$ commutes with $A$, that is, $AX=XA$, if and only if there exist two complex numbers $\alpha,\alpha'$ such that $X=\alpha A+\alpha' I_2$.
\item Let $A,B$ and $C$ be matrices from $M_2(\mathbb{C})$, such that $AB\ne BA$, $AC=CA$ and $BC=CB$. Prove that $C$ commutes with all matrices from $M_2(\mathbb{C})$.
\end{enumerate}
\end{problem}

\begin{problem}
Let $A,B,C\in M_3(\mathbb{R})$ such that $\det(A)=\det(B)=\det(C)$ and $\det(A+iB)=\det(C+iA)$. Prove that $\det(A+B)=\det(C+A)$.
\end{problem}

\begin{problem}
Let $A\in M_n(\mathbb{R})$. If $A^TA=I_n$, prove that: a) $|\operatorname{Tr}(A)|\le n$; b) If $n$ is odd, then $\det(A^2-I_n)=0$.
\end{problem}

\begin{problem}
Consider the matrices $A,B\in M_3(\mathbb{C})$ with $A=-A^T$ and $B=B^T$. Prove that if the polynomial function defined by $f(x)=\det(A+xB)$ has a multiple root, then $\det(A+B)=\det B$.
\end{problem}

\begin{problem}
Let $A,B\in M_2(\mathbb{C})$ be two non-zero matrices such that $AB+BA=O_2$ and $\det(A+B)=0$. Prove that $A$ and $B$ have null traces.
\end{problem}

\begin{problem}
Let $A,B$ be $2\times 2$ matrices with real elements such that $AB=A^2B^2-(AB)^2$ and $\det(B)=2$.
\begin{enumerate}[label=(\alph*)]
\item Prove that the matrix $A$ is not invertible.
\item Calculate $\det(A+2B) - \det(B+2A)$.
\end{enumerate}
\end{problem}

\begin{problem}
\begin{enumerate}[label=(\alph*)]
\item Give an example of matrices $A$ and $B$ from $M_2(\mathbb{R})$ such that $A^2+B^2 = \begin{pmatrix}2 & 3\\ 3 & 2\end{pmatrix}$.
\item Let $A$ and $B$ be matrices from $M_2(\mathbb{R})$ such that $A^2+B^2=\begin{pmatrix}2 & 3\\ 3 & 2\end{pmatrix}$. Show that $AB\ne BA$.
\end{enumerate}
\end{problem}

\begin{problem}
Show that the function $f:M_2(\mathbb{C})\to M_2(\mathbb{C})$, $f(X)=X^n$, is not injective nor surjective for all $n\ge 2$.
\end{problem}

\begin{problem}
Let $A,B\in M_2(\mathbb{R})$. Prove that $\det(A^2+B^2) \ge \det(AB-BA)$.
\end{problem}

\begin{problem}
Let $A,B\in M_2(\mathbb{R})$ such that $\det(AB+BA)\le 0$. Show that $\det(A^2+B^2)\ge 0$.
\end{problem}

\begin{problem}
Let $A\in M_3(\mathbb{R})$ such that $\operatorname{r}(A)=\operatorname{tr}(A^2)=0$. Show that $\det(A^2+I_3) = (\det(A))^2+1$.
\end{problem}

\begin{problem}
Let $A\in M_3(\mathbb{C})$ such that $\operatorname{tr}(A)=\operatorname{tr}(A^2)=0$. Show that if $|\det(A)|<1$ then the matrix $I_3+A^n$ is invertible for all $n\ge 1$.
\end{problem}

\begin{problem}
Let $A,B\in M_2(\mathbb{R})$ be two matrices with positive entries. Show that $(AB)^2=(BA)^2$ if and only if $AB=BA$.
\end{problem}

\begin{problem}
Let $A\in M_2(\mathbb{R})$ such that $\det(A)\ge 0$. Show that
\[
\det(A^n+I_2) + \det(A^n-I_2) \ge \frac{1}{2^{n-2}}(\det(A)+1)^n.
\]
\end{problem}

\begin{problem}
Let $A\in M_2(\mathbb{R})$ be a matrix such that $A\ne\operatorname{adj}(A)$ and $A^3 = (\operatorname{adj}(A))^3$. Show that $\det(A)=\operatorname{tr}^2(A)$.
\end{problem}

\begin{problem}
Find all polynomials $P$ with real coefficients with the following property: for any two matrices $A,B\in M_2(\mathbb{R})$ with $A\ne B$ we have $P(A)\ne P(B)$.
\end{problem}

\begin{problem}
Let $A,B\in M_2(\mathbb{Z})$ with $AB=BA$ such that $\det(A)=\det(B)=0$. Show that $\det(A^3+B^3)$ is the cube of an integer.
\end{problem}

\begin{problem}
Let $n\ge 3$ be a positive integer. Determine $X\in M_2(\mathbb{R})$ such that
\[
X^n+X^{n-2} = \begin{pmatrix}1 & -1\\ -1 & 1\end{pmatrix}.
\]
\end{problem}

\begin{problem}
Let $A\in M_2(\mathbb{R})$ such that $\operatorname{tr}(A)>2$. Show that $A^n\ne I_2$ for all $n\ge 1$.
\end{problem}

\begin{problem}
\begin{enumerate}[label=(\alph*)]
\item Let $A\in M_3(\mathbb{C})$ such that $A^2=O_3$. Show that $\operatorname{tr}(A)=0$.
\item Let $A\in M_3(\mathbb{C})$ such that $\operatorname{tr}(A)=\operatorname{tr}(A^2)=0$. Show that $\det(A)=0$.
\end{enumerate}
\end{problem}

\begin{problem}
Show that the equation $X^n=I_2$ has at least $n$ solutions in $M_2(\mathbb{R})$ with $n\ge 1$.
\end{problem}

\begin{problem}
Find all matrices $A\in M_2(\mathbb{Z})$ such that $\det(A^3+I_2)=1$.
\end{problem}

\begin{problem}
Let $A,B\in M_2(\mathbb{Z})$ such that $AB=BA$ and $\det(B)=1$. Show that if $\det(A^3+B^3)=1$, then $A^2=O_2$.
\end{problem}

\begin{problem}
Let $A,B\in M_2(\mathbb{R})$ with $AB=BA$. If there exists a positive integer $n$ such that $\det(A^{2n}+B^{2n})=0$ then the matrices $A$ and $B$ are simultaneously invertible or non-invertible.
\end{problem}

\begin{problem}
Let $A,B\in M_2(\mathbb{R})$ such that $A^2+B^2=2AB$. Prove that $AB=BA$ and $\operatorname{tr}(A)=\operatorname{tr}(B)$.
\end{problem}

\begin{problem}
Show that if $A,B,C\in M_3(\mathbb{C})$ verify $A^3=B^3=C^3=O_3$ and they commute pairwise, then $ABC=O_3$.
\end{problem}

\begin{problem}
Let $A,B\in M_3(\mathbb{C})$. Show that $\det(AB-BA) = \operatorname{tr}\big(AB(BA-AB)BA\big)$.
\end{problem}

\begin{problem}
Let $A,B\in M_n(\mathbb{R})$, $A\ne B$, such that $A^3=B^3$ and $A^2B=AB^2$. Is it possible that the matrix $A^2+B^2$ is invertible?
\end{problem}

\begin{problem}
Let $A,B\in M_n(\mathbb{R})$ such that $(AB)^2=O_n$. Is it true that $(BA)^2=O_n$?
\end{problem}

\begin{problem}
Let $A,B\in M_n(\mathbb{R})$ such that $AA^T=I_n$ and $BB^T=I_n$, with $n$ odd. Show that at least one of the matrices $A+B$ and $A-B$ is singular.
\end{problem}

\begin{problem}
Let $A,B\in M_n(\mathbb{R})$ with $B^2=O_n$. Show that $\det(AB+BA+I_n) \ge 0$.
\end{problem}

\begin{problem}
\begin{enumerate}[label=(\alph*)]
\item Find two matrices $A,B\in M_2(\mathbb{C})$ such that $A^2+B^2=I_2$ and $AB-BA$ is nonsingular.
\item Let $n$ be odd and $A,B\in M_n(\mathbb{C})$ such that $A^2+B^2=I_n$. Show that $\det(AB-BA)\ge 0$.
\end{enumerate}
\end{problem}

\begin{problem}
Let $A,B,C\in M_2(\mathbb{R})$ be matrices that commute pairwise such that $\det(A)=0$. Show that $\det(A^2+B^2+C^2)\ge 0$.
\end{problem}

\begin{problem}
Show that if $A\in M_3(\mathbb{Q})$ such that $A^8=I_3$, then $A^4=I_3$.
\end{problem}

\begin{problem}
Let $A\in M_n(\mathbb{R})$ such that $|\operatorname{tr}(A)|>n$. Show that $A^k\ne A^p$ for all $k,p\in\mathbb{N}$, $k\ne p$.
\end{problem}

\begin{problem}
Let $A$ be an $n\times n$ matrix with real entries. Show that $\operatorname{tr}(A^k)=0$ for all $k=0,1,\dots,n$ if and only if $A^n=O_n$.
\end{problem}

\begin{problem}
\begin{enumerate}[label=(\alph*)]
\item Are there matrices $A,B\in M_3(\mathbb{C})$ such that $(AB-BA)^2=I_3$?
\item Are there matrices $A,B\in M_3(\mathbb{C})$ such that $(AB-BA)^{1993}=I_3$?
\end{enumerate}
\end{problem}

\begin{problem}
Let $A$ and $B$ be two matrices from $M_3(\mathbb{C})$ such that $(AB)^2=A^2B^2$ and $(BA)^2=B^2A^2$. Prove that $(AB-BA)^3=O_3$.
\end{problem}

\begin{problem}
Let $A,B$ be two matrices from $M_3(\mathbb{C})$ such that $A^2=AB+BA$. Prove that the matrix $AB-BA$ is singular.
\end{problem}

\begin{problem}
Let $A,B\in M_k(\mathbb{R})$ with $AB=BA$. Show that $\det(A+B)\ge 0$ if and only if $\det(A^n+B^n)\ge 0$ for all $n\ge 1$.
\end{problem}

\begin{problem}
Let $r,s$ be odd prime numbers and $A,B\in M_n(\mathbb{C})$ such that $AB=BA$, $A^r=I_n$ and $B^s=I_n$. Show that the matrix $A+B$ is invertible.
\end{problem}

\begin{problem}
Let $n,k\ge 1$ and $A_1,A_2,\dots,A_k\in M_n(\mathbb{R})$. Show that $\det\!\left(\sum_{i=1}^k A_i^TA_i\right)\ge 0$.
\end{problem}

\begin{problem}
Let $A,B,C,D\in M_n(\mathbb{C})$ where $A$ and $C$ are invertible. If $A^kB=C^kD$ for all $k\ge 1$ then show that $B=D$.
\end{problem}

\begin{problem}
Let $A_0,A_1,\dots,A_n\in M_2(\mathbb{R})$, $n\ge 2$, be $n+1$ matrices such that $A_0\ne a\cdot A_2$ for any real $a$ and $A_0A_k=A_kA_0$ for all $k=1,2,\dots,n$. Show that
\begin{enumerate}[label=(\alph*)]
\item $\det\!\left(\sum_{k=1}^n A_k^2\right) \ge 0$;
\item If $\det\!\left(\sum_{k=1}^n A_k^2\right)=0$ and $A_2\ne a\cdot A_1$ for any real number $a$, then $\sum_{k=1}^n A_k^2 = O_n$.
\end{enumerate}
\end{problem}

\begin{problem}
Let $A\in M_2(\mathbb{C})$ and $C(A)=\{B\in M_2(\mathbb{C}): AB=BA\}$. Show that $|\det(A+B)| \ge |\det(B)|$ for any $B\in C(A)$ if and only if $A^2=O_2$.
\end{problem}

\begin{problem}
Let $M = \{A\in M_2(\mathbb{C}) : \det(A-zI_2)=0 \Rightarrow |z|<1\}$. Show that if $A,B\in M$ and $AB=BA$, then $AB\in M$.
\end{problem}

\begin{problem}
Let $A\in M_n(\mathbb{Z})$. Show that $\det(A^4+I_n) \ne 13$.
\end{problem}

\begin{problem}
Let $A,B\in M_n(\mathbb{C})$. Show that if the equation $AX+XB=C$, with $X,C\in M_n(\mathbb{C})$, has a solution for all $C\in M_n(\mathbb{C})$, then $\det(A)\ne 0$ or $\det(B)\ne 0$.
\end{problem}

\begin{problem}
Let $A\in M_n(\mathbb{Z})$, $m>1$, such that $\gcd(m,\det(A))=1$. Show that the matrix $A^2+A+m\cdot I_n$ is nonsingular.
\end{problem}

\begin{problem}
Let $A\in M_2(\mathbb{Z})$ different from the null matrix. Show that the following assertions are equivalent: (i) $4\det(A)=(\operatorname{tr}(A))^2$ and $A\ne k\cdot I_2$ for $k\in\mathbb{Z}$; (ii) for any matrix $B\in M_2(\mathbb{Z})$ which commutes with $A$, $\det(A^2+B^2)$ is the square of an integer.
\end{problem}

\begin{problem}
Consider the matrices $A,B\in M_n(\mathbb{C})$ and the polynomial $f\in\mathbb{C}[X]$, $f(X)=\det(A+X\cdot B)$. Show that the degree of the polynomial $f$ is less than or equal to the rank of the matrix $B$.
\end{problem}

\begin{problem}
\begin{enumerate}[label=(\alph*)]
\item Let $A$ be a square matrix of size $3$, with real entries. Show that if $f$ is a polynomial with real coefficients, but with no real roots, then $f(A)\ne O_3$.
\item Show that there exists a positive integer $n$ such that $(A-\operatorname{adj}(A))^{2n} + A^{2n} + (\operatorname{adj}(A))^{2n} = O_3$ if and only if $\det(A)=0$.
\end{enumerate}
\end{problem}

\begin{problem}
\begin{enumerate}[label=(\alph*)]
\item Show that any matrix $A\in M_4(\mathbb{C})$ can be written as the sum of four matrices $B_i\in M_4(\mathbb{C})$, $i=1,2,3,4$, of rank $1$.
\item Show that $I_4$ cannot be written as the sum of four matrices of rank $1$.
\end{enumerate}
\end{problem}

\begin{problem}
Let $A,B\in M_n(\mathbb{C})$ such that there exists $k>1$ with $B^k=O_n$. Show that if $AB=BA$, then $\det(A+B)=\det(A)$.
\end{problem}

\begin{problem}
Let $A,B$ be two $n\times n$ matrices with complex entries, and $p,q$ fixed positive integers such that $AB=BA$ and $A^p=B^q=O_n$. Show that the matrix $A+B+I_n$ is invertible.
\end{problem}

\begin{problem}
Let $m,n>2$ be integers. The matrices $A_1,A_2,\dots,A_m\in M_n(\mathbb{R})$ are not all nilpotent. Prove that there is an integer $k>0$ such that $A_1^k+A_2^k+\cdots+A_m^k \ne O_n$.
\end{problem}

\begin{problem}
Let $A,B$ be $3\times 3$ matrices with real entries such that $AB=O_3$.
\begin{enumerate}[label=(\alph*)]
\item Prove that $f:\mathbb{C}\to\mathbb{C}$, $f(x)=\det(A^2+B^2+x\cdot BA)$ is a polynomial function of degree not greater than $2$.
\item Prove that $\det(A^2+B^2)\ge 0$.
\end{enumerate}
\end{problem}

\begin{problem}
Consider a positive integer $n$ and $A,B$ two matrices in $M_n(\mathbb{C})$ such that $A^2+B^2=2AB$. Prove that:
\begin{enumerate}[label=(\alph*)]
\item The additive commutator $AB-BA$ is not invertible;
\item If the matrix $A-B$ has rank $1$, then matrices $A$ and $B$ commute.
\end{enumerate}
\end{problem}

\begin{problem}
Let $A$ be an invertible matrix in $M_4(\mathbb{R})$ such that $\operatorname{tr}(A)=\operatorname{tr}(\operatorname{adj}(A))\ne 0$. Prove that the matrix $A^2+I_4$ is singular if and only if there exists a nonzero matrix $B$ in $M_4(\mathbb{R})$, so that $AB=-BA$.
\end{problem}

\begin{problem}
Consider two nonnegative integers $n$ and $k$ such that $n\ge 2$ and $1\le k\le n-1$. Suppose that a matrix $A\in M_n(\mathbb{C})$ has exactly $k$ minors of order $n-1$ which are null. Prove that $\det(A)\ne 0$.
\end{problem}

\begin{problem}
Let $A$ be a non-invertible square matrix of size $n$, $n>2$, with real entries, and let $\operatorname{adj}(A)$ be the adjugate of $A$. Prove that $\operatorname{tr}(\operatorname{adj}(A)) \ne -1$ if and only if the matrix $I_n+\operatorname{adj}(A)$ is invertible.
\end{problem}

\begin{problem}
Consider $1\le a_1<a_2<a_3<a_4$ and $x_1<x_2<x_3<x_4$ be real numbers and $M = (a_i^{x_j})_{1\le i,j\le 4}$ be a square matrix of size $4$. Show that $\det(M)>0$.
\end{problem}

\begin{problem}
Let $S(X)$ be the sum of all the entries of the matrix $X\in M_n(\mathbb{R})$. Show that for any matrices $A_1,A_2,\dots,A_k\in M_n(\mathbb{R})$ we have the following inequality:
\[
S^2(A_1A_2^T)\cdot S^2(A_2A_3^T) \cdots S^2(A_kA_1^T) \le S^2(A_1A_1^T)\cdot S^2(A_2A_2^T)\cdots S^2(A_kA_k^T).
\]
\end{problem}

\begin{problem}
Let $n>k$ and let $A_1,A_2,\dots,A_k$ be real $n\times n$ matrices of rank $n-1$. Prove that $A_1A_2\cdots A_k \ne$ (has rank strictly less than $n$; state and prove the appropriate rank bound).
\end{problem}

\begin{problem}
Two matrices $A,B$ of the same size having real entries satisfy $A^{2002}=B^{2003}=I_n$ and $AB=BA$. Show that $A+B+I_n$ is invertible.
\end{problem}

\begin{problem}
Let $A$ be an $n\times n$ matrix of real numbers for some $n\ge 1$. For each positive integer $k$, let $A^{[k]}$ be the matrix obtained by raising each entry to the $k$-th power. Show that if $A^k=A^{[k]}$ for $k=1,2,\dots,n+1$, then $A^k=A^{[k]}$ for all $k\ge 1$.
\end{problem}

\begin{problem}
Let $A$ and $B$ be $2\times 2$ matrices with integer entries such that $A,A+B,A+2B,A+3B$, and $A+4B$ are all invertible matrices whose inverses have integer entries. Show that $A+5B$ is invertible and that its inverse has integer entries.
\end{problem}

\begin{problem}
Let $A,B,C\in M_n(\mathbb{R})$ such that $ABC=O_n$ and $\operatorname{rank}(B)=1$. Prove that $AB=O_n$ or $BC=O_n$.
\end{problem}

\begin{problem}
Let $A,B\in M_n(\mathbb{C})$ such that $AB=BA$ and $\det(B)\ne 0$.
\begin{enumerate}[label=(\alph*)]
\item If $|\det(A+zB)|=1$ for any $z\in\mathbb{C}$ such that $|z|=1$ then $A^n=O_n$.
\item Is the statement from (a) still true if $AB\ne BA$?
\end{enumerate}
\end{problem}

\begin{problem}
Let $A$ be an $n\times n$ matrix with complex elements. Prove that $A^{-1}=A$ if and only if there exists an invertible matrix $B$ with complex elements such that $A=B^{-1}\overline{B}$.
\end{problem}

\begin{problem}
Let $A=(a_{ij})_{1\le i,j\le n}$ be a real $n\times n$ matrix, such that $a_{ij}+a_{ji}=0$ for all $i,j$. Prove that for all non-negative real numbers $x,y$ we have
\[
\det(A+xI_n)\cdot\det(A+yI_n) \ge \det(A+\sqrt{xy}\,I_n)^2.
\]
\end{problem}

\begin{problem}
Let $A,B\in M_2(\mathbb{R})$ that satisfy $A^2+B^2=AB$. Prove that $(AB-BA)^2 = O_2$.
\end{problem}

\begin{problem}
Let $n\in\mathbb{N}$, $n\ge 2$.
\begin{enumerate}[label=(\alph*)]
\item Give an example of two matrices $A,B\in M_n(\mathbb{C})$ such that $\operatorname{rank}(AB)-\operatorname{rank}(BA) = \lfloor n/2\rfloor$.
\item Prove that for all matrices $X,Y\in M_n(\mathbb{C})$ we have $\operatorname{rank}(XY) - \operatorname{rank}(YX) \le \lfloor n/2\rfloor$.
\end{enumerate}
\end{problem}

\begin{problem}
\begin{enumerate}[label=(\alph*)]
\item Find a matrix $A\in M_3(\mathbb{C})$ such that $A^2\ne O_3$ and $A^3=O_3$.
\item Let $n,p\in\{2,3\}$. Prove that if there is a bijective function $f:M_n(\mathbb{C})\to M_p(\mathbb{C})$ such that $f(XY)=f(X)f(Y)$ for all $X,Y\in M_n(\mathbb{C})$, then $n=p$.
\end{enumerate}
\end{problem}

\begin{problem}
\begin{enumerate}[label=(\alph*)]
\item Let $A,B\in M_3(\mathbb{R})$ such that $\operatorname{rank}(A) > \operatorname{rank}(B)$. Prove that $\operatorname{rank}(A^2) \ge \operatorname{rank}(B^2)$.
\item Find the non-constant polynomials $f\in\mathbb{R}[X]$ such that for all $A,B\in M_4(\mathbb{R})$ with $\operatorname{rank}(A)>\operatorname{rank}(B)$, we have $\operatorname{rank}(f(A)) > \operatorname{rank}(f(B))$.
\end{enumerate}
\end{problem}

\begin{problem}
Let $p\ge 2$ be an integer and $A$ an $n\times n$ matrix with real elements such that $A^{p+1}=A$.
\begin{enumerate}[label=(\alph*)]
\item Prove that $\operatorname{rank}(A) + \operatorname{rank}(I_n-A^p) = n$.
\item Prove that if $p$ is prime then $\operatorname{rank}(I_n-A) = \operatorname{rank}(I_n-A^2) = \cdots = \operatorname{rank}(I_n-A^{p-1})$.
\end{enumerate}
\end{problem}

\begin{problem}
Let $A,B\in M_n(\mathbb{R})$. Prove that $\operatorname{rank}(A)+\operatorname{rank}(B) < n$ if and only if there exists an invertible matrix $X\in M_n(\mathbb{R})$ such that $AXB=O_n$.
\end{problem}

\begin{problem}
Let $S_1,S_2,\dots,S_{2^n-1}$ be the nonempty subsets of $\{1,2,\dots,n\}$ in some order, and let $M$ be the $(2^n-1)\times(2^n-1)$ matrix whose $(i,j)$ entry is
\[
m_{ij} = \begin{cases} 0 & \text{if } S_i\cap S_j = \emptyset\\ 1 & \text{otherwise.}\end{cases}
\]
Determine $\det(M)$.
\end{problem}

\begin{problem}
Let $n$ be a positive integer. Suppose that $A,B,M$ are $n\times n$ matrices with real entries such that $AM=MB$ and such that $A$ and $B$ have the same characteristic polynomial. Prove that $\det(A-MX) = \det(B-XM)$ for every $n\times n$ matrix $X$ with real entries.
\end{problem}

\begin{problem}
Let $S$ be the set of all $2\times 2$ matrices $M = \begin{pmatrix}a & b\\ c & d\end{pmatrix}$ whose entries $a,b,c,d$ (in that order) form an arithmetic progression. Find all matrices $M$ in $S$ for which there is some integer $k>1$ such that $M^k$ is also in $S$.
\end{problem}

\begin{problem}
Let $A$ be the $n\times n$ matrix whose entry in the $i$-th row and $j$-th column is $\dfrac{1}{\min(i,j)}$ for $1\le i,j\le n$. Compute $\det(A)$.
\end{problem}

\begin{problem}
Let $A$ be an $m\times n$ matrix with rational entries. Suppose that there are at least $m+n$ distinct prime numbers among the absolute values of the entries of $A$. Show that the rank of $A$ is at least $2$.
\end{problem}

\begin{problem}
Let $d_n$ be the determinant of the $n\times n$ matrix whose entries, from left to right and then from top to bottom, are $\cos 1,\cos 2,\dots,\cos n^2$. (For example,
\[
d_3 = \begin{vmatrix}\cos1 & \cos2 & \cos3\\ \cos4 & \cos5 & \cos6\\ \cos7 & \cos8 & \cos9\end{vmatrix},
\]
where the argument of $\cos$ is always in radians, not degrees.) Evaluate $\lim_{n\to\infty} d_n$.
\end{problem}

\begin{problem}
Let $Z$ denote the set of points in $\mathbb{R}^n$ whose coordinates are $0$ or $1$. (Thus $Z$ has $2^n$ elements, which are vertices of a hypercube in $\mathbb{R}^n$.) Given a vector subspace $V$ of $\mathbb{R}^n$, let $Z(V)$ denote the number of members of $Z$ that lie in $V$. Let $k$ be given, $0<k<n$. Find the maximum, over all vector subspaces $V\subset\mathbb{R}^n$ of dimension $k$, of the number of points in $V\cap Z$.
\end{problem}

\begin{problem}
Let $H$ be an $n\times n$ matrix all of whose entries are $\pm 1$ and whose rows are mutually orthogonal. Suppose $H$ has an $a\times b$ submatrix whose entries are all $1$. Show that $ab\le n$.
\end{problem}

\begin{problem}
Define a sequence $\{u_n\}_{n=0}^\infty$ by $u_0=u_1=u_2=1$, and thereafter by the condition that
\[
\det\begin{pmatrix}u_n & u_{n+1}\\ u_{n+2} & u_{n+3}\end{pmatrix} = n!
\]
for all $n\ge 0$. Show that $u_n$ is an integer for all $n$. (By convention, $0!=1$.)
\end{problem}

\begin{problem}
For an integer $n\ge 3$ let $\theta=2\pi/n$. Evaluate the determinant of the $n\times n$ matrix $I+A$, where $I$ is the $n\times n$ identity matrix and $A=(a_{jk})$ has entries $a_{jk} = \cos(j\theta+k\theta)$ for all $j,k$.
\end{problem}

\begin{problem}
For any square matrix $A$, we can define $\sin A$ by the usual power series $\sum_{n=0}^\infty \dfrac{(-1)^n}{(2n+1)!}A^{2n+1}$. Prove or disprove: there exists a $2\times 2$ matrix $A$ with real entries such that
\[
\sin A = \begin{pmatrix}1 & 1996\\ 0 & 1\end{pmatrix}.
\]
\end{problem}

\begin{problem}
Let $D_n$ denote the value of the $(n-1)\times(n-1)$ determinant
\[
D_n = \begin{vmatrix} 1 & 1 & 1 & \cdots\\ 1 & 1 & 1 & \cdots\\ 1 & 1 & 1 & \cdots\\ \vdots & \vdots & \vdots & \ddots\end{vmatrix}.
\]
Is the set $\{D_n/n! : n\ge 2\}$ bounded?
\end{problem}

\begin{problem}
Let $M$ be a set of real $n\times n$ matrices such that (i) $I\in M$, where $I$ is the $n\times n$ identity matrix; (ii) if $A\in M$ and $B\in M$ then either $AB\in M$ or $-AB\in M$ but not both; (iii) if $A\in M$ and $B\in M$ then either $AB=BA$ or $AB=-BA$; (iv) if $A\in M$ and $A\ne I$ then there is at least one $B\in M$ such that $AB=-BA$. Prove that $M$ contains at most $n^2$ matrices.
\end{problem}

\begin{problem}
Let $A$ and $B$ be different $n\times n$ matrices with real entries. If $A^3=B^3$ and $A^2B=B^2A$, can $A^2+B^2$ be invertible?
\end{problem}

\begin{problem}
If $A$ and $B$ are square matrices of the same size such that $ABAB=O$, does it follow that $BABA=O$?
\end{problem}

\begin{problem}
A transversal of an $n\times n$ matrix $A$ consists of $n$ entries of $A$, no two in the same row or column. Let $f(n)$ be the number of $n\times n$ matrices $A$ satisfying the following two conditions: (a) each entry $a_{i,j}$ of $A$ is in the set $\{-1,0,1\}$; (b) the sum of the $n$ entries of a transversal is the same for all transversals of $A$. An example of such a matrix $A$ is
\[
A = \begin{pmatrix}-1 & 0 & -1\\ 0 & 1 & 0\\ 0 & 1 & 0\end{pmatrix}.
\]
Determine with proof a formula for $f(n)$ of the form $f(n) = a_1b_1^n + a_2b_2^n + a_3b_3^n + a_4$ where the $a_i$ and $b_i$ are rational numbers.
\end{problem}

\begin{problem}
Let $G$ be a finite set of real $n\times n$ matrices $\{M_i\}$, $1\le i\le r$, which form a group under matrix multiplication. Suppose that $\sum_{i=1}^r \operatorname{tr}(M_i) = 0$ where $\operatorname{tr}(A)$ denotes the trace of the matrix $A$. Prove that $\sum_{i=1}^r M_i$ is the $n\times n$ zero matrix.
\end{problem}

\begin{problem}
Let $k \in \mathbb{N}^*$, $u_1, u_2, \ldots, u_k \in \mathbb{C}^*$, $z_1, z_2, \ldots, z_k \in \mathbb{C}^*$ be distinct, and let
$$a_n = u_1 z_1^n + u_2 z_2^n + \cdots + u_k z_k^n, \quad n \in \mathbb{N}.$$
Show that if the set $\{a_n \mid n \in \mathbb{N}\}$ is finite, then there exists $d \in \mathbb{N}^*$ such that $a_n = a_{n+d}$ for every $n \in \mathbb{N}$.
\end{problem}

\begin{problem}
Let $A \in M_n(\mathbb{R})$ with $\det A \neq 0$, and let $f_A : M_n(\mathbb{R}) \to \mathbb{R}$,
$$f_A(X) = \sum_{i,j=1}^{n} b_{ij} x_{ij}, \quad \text{where } B = AA^t.$$
Show that:
\begin{enumerate}[label=(\alph*)]
\item $f_A(XX^t) = 0 \Rightarrow X = 0$,
\item $[f_A(XY)]^2 \le f_A(XX^t)\, f_A(Y^tY)$, for all $X, Y \in M_n(\mathbb{R})$.
\end{enumerate}
\end{problem}

\begin{problem}
Consider the function $f : M_{n,1}(\mathbb{R}) \to [0,\infty)$,
$$f(X) = \sum_{k=1}^{n} |x_k|,$$
for every $X \in M_{n,1}(\mathbb{R})$. Show that there is only a finite number of matrices $A \in M_n(\mathbb{R})$ with the property $f(AX) = f(X)$ for every $X \in M_{n,1}(\mathbb{R})$.
\end{problem}

\begin{problem}
Consider the matrix
$$A_n = \begin{pmatrix}
1 & \tfrac{1}{2} & \tfrac{1}{3} & \cdots & \tfrac{1}{n} \\[4pt]
\tfrac{1}{2} & \tfrac{1}{3} & \tfrac{1}{4} & \cdots & \tfrac{1}{n+1} \\[4pt]
\cdots & \cdots & \cdots & \cdots & \cdots \\[4pt]
\tfrac{1}{n} & \tfrac{1}{n+1} & \tfrac{1}{n+2} & \cdots & \tfrac{1}{2n-1}
\end{pmatrix},$$
where $n$ is a nonzero natural number.
\begin{enumerate}[label=\arabic*)]
\item Show that $\det A_n$ is nonzero.
\item Show that the sum of the entries of the inverse matrix $A_n^{-1}$ equals $n^2$.
\end{enumerate}
\end{problem}

\begin{problem}
Let $A, B \in M_n(\mathbb{C})$ and $\varepsilon = \cos\dfrac{2\pi}{n} + i \sin \dfrac{2\pi}{n}$. Show that
$$\det(A+B) + \det(A+\varepsilon B) + \cdots + \det(A + \varepsilon^{n-1} B) = n\left[\det A + \det B\right].$$
\end{problem}

\begin{problem}
Show that for every matrix $A \in M_n(\mathbb{C})$ there exist infinitely many natural numbers $k$ such that the matrix $A^k + I_n$ is invertible.
\end{problem}

\begin{problem}
Let $A, B \in M_n(\mathbb{C})$ have the property
$$|\det(A + zB)| \le 1 \quad \text{for every } z \in \mathbb{C} \text{ with } |z| = 1.$$
Show that:
\begin{enumerate}[label=(\alph*)]
\item $|\det A + z \det B| \le 1$ for every $z \in \mathbb{C}$ with $|z| = 1$.
\item If $A, B \in M_n(\mathbb{R})$, then $(\det A)^2 + (\det B)^2 \le 1$.
\end{enumerate}
\end{problem}

\begin{problem}
Let $A \in M_n(\mathbb{R})$ be a matrix with positive entries satisfying
$$\sum_{k=1}^{n} a_{ik} = 1, \quad i = \overline{1,n}.$$
Show that $A$ cannot have eigenvalues of modulus greater than $1$.
\end{problem}

\begin{problem}
Let the matrices $A, B, C \in M_n(\mathbb{C})$ be such that $C$ commutes with $A$ or with $B$. Show that
$$\det(AB + C) = \det(BA + C).$$
\end{problem}

\begin{problem}
Let $A \in M_n(\mathbb{Q})$ be such that $\det(A - I_n) = 0$ and $A^p = I_n$, where $p$ is a prime number. Show that $p - 1 \mid n$.
\end{problem}

\begin{problem}
Let $S(X)$ denote the sum of all entries of the matrix $X$. Show that if $A \in M_n(\mathbb{C})$ and
\begin{enumerate}[label=(\alph*)]
\item $S(A) = S(A^2) = \cdots = S(A^n) = n$, then $S(A^k) = n$ for every $k \in \mathbb{N}^*$.
\item $S(A) = n^2,\ S(A^2) = n^3,\ \ldots,\ S(A^n) = n^{n+1}$, then $S(A^k) = n^{k+1}$ for every $k \in \mathbb{N}^*$.
\end{enumerate}
\end{problem}

\begin{problem}
Show that if there exists an invertible matrix $A \in M_n(\mathbb{Q})$ such that
$$A^{-1} = A^2 + A,$$
then $n$ is divisible by $3$.
\end{problem}

\begin{problem}
Let $\displaystyle\sum_{n=0}^{\infty} a_n x^{2n}$ be the series expansion of the function $f(x) = \dfrac{1}{\cos x}$. The numbers $e_n = \dfrac{1}{(2n)!} a_n$ are called the \emph{Euler numbers}. Show that $e_n = (2n)!\, D_n$, where
$$D_n = \begin{vmatrix}
\dfrac{1}{2!} & 1 & 0 & 0 & \cdots & 0 \\[6pt]
\dfrac{1}{4!} & \dfrac{1}{2!} & 1 & 0 & \cdots & 0 \\[6pt]
\dfrac{1}{6!} & \dfrac{1}{4!} & \dfrac{1}{2!} & 1 & \cdots & 0 \\[6pt]
\cdots & \cdots & \cdots & \cdots & \cdots & \cdots \\[6pt]
\dfrac{1}{(2n)!} & \dfrac{1}{(2n-2)!} & \dfrac{1}{(2n-4)!} & \dfrac{1}{(2n-6)!} & \cdots & \dfrac{1}{2!}
\end{vmatrix}.$$
\end{problem}

\begin{problem}
Let $n \ge 2$, $n \in \mathbb{N}$. Determine all polynomials $f \in \mathbb{C}[X]$ with the property
\begin{enumerate}[label=(\alph*)]
\item $f(\operatorname{Tr} A) = \operatorname{Tr}(f(A))$ for every $A \in M_n(\mathbb{C})$;
\item $\det(f(A)) = f(\det A)$ for every $A \in M_n(\mathbb{C})$.
\end{enumerate}
\end{problem}

\begin{problem}
Determine the number of $(m,n)$-type matrices with entries $a_{ij} \in \{\pm 1\}$ for which the product of the entries of each row and of each column equals $-1$.
\end{problem}

\begin{problem}
Let $A \in M_3(\mathbb{C})$ with $\det A = 1$, $\operatorname{Tr} A = \operatorname{Tr} A^{-1} = 0$. Show that $A^3 = I_3$.
\end{problem}

\begin{problem}
Let $A, B \in M_n(\mathbb{C})$ be such that $AB + A + B = 0$. Show that $AB = BA$.
\end{problem}

\begin{problem}
Let $A, B \in M_n(\mathbb{R})$ be such that $A^2 + B^2 = AB$ and the matrix $AB - BA$ is invertible. Show that $n$ is divisible by $3$.
\end{problem}

\begin{problem}
Let $A, B, C \in M_n(\mathbb{C})$ with $A$ invertible and $(A-B)C = BA^{-1}$. Show that $C(A-B) = A^{-1}B$.
\end{problem}

\begin{problem}
Let $A_1, A_2, \ldots, A_k \in M_n(\mathbb{R})$ be such that
$$A_1 A_1^t + A_2 A_2^t + \cdots + A_n A_n^t = 0.$$
Show that $A_1 = A_2 = \cdots = A_n = 0$.
\end{problem}

\begin{problem}
Let $A, B, C, D \in M_2(\mathbb{C})$ and denote $[A,B] = AB - BA$. Show that there exists $\lambda \in \mathbb{C}$ such that
$$[A,B]\cdot[C,D] - [C,D]\cdot[A,B] = \lambda I_2.$$
\end{problem}

\begin{problem}
Let $A \in M_{3,2}(\mathbb{C})$, $B \in M_{2,3}(\mathbb{C})$ be such that
$$AB = \begin{pmatrix} 8 & 2 & -2 \\ 2 & 5 & 4 \\ -2 & 4 & 5 \end{pmatrix}.$$
Show that $BA = 9 I_2$.
\end{problem}

\begin{problem}
Let $A \in M_n(\mathbb{R})$, $A \neq 0$, be such that $a_{ik} a_{jk} = a_{kk} a_{ij}$ for all $i,j,k = \overline{1,n}$. Show that:
\begin{enumerate}[label=(\alph*)]
\item $\operatorname{Tr} A \neq 0$.
\item $A^t = A$.
\item $f_A(x) = x^{n-1}(x - \operatorname{Tr} A)$.
\end{enumerate}
\end{problem}

\begin{problem}
Let $A, B \in M_n(\mathbb{R})$ be such that $\operatorname{Tr}(AA^t + BB^t) = \operatorname{Tr}(AB + A^tB^t)$. Show that $A = B^t$.
\end{problem}

\begin{problem}
Let $A, B \in M_n(\mathbb{C})$. Show that
$$\operatorname{Tr}(AB) \le \frac{1}{2}\operatorname{Tr}(AA^* + BB^*).$$
\end{problem}

\begin{problem}
Let $A \in M_n(\mathbb{C})$ be a nonsingular matrix with columns $A_1, A_2, \ldots, A_n$, and let $B \in M_n(\mathbb{C})$ have columns $A_2, A_3, \ldots, A_n, 0$. Show that the matrices $C = BA^{-1}$ and $D = A^{-1}B$ both have rank $n-1$ and all of their eigenvalues equal to zero.
\end{problem}

\begin{problem}
Let $\lambda_1,\lambda_2,\dots,\lambda_n$ and $a,b$ be real parameters. Consider the polynomial
\[
p(x) = \det\begin{pmatrix}
\lambda_1+x & a+x & \cdots & a+x\\
b+x & \lambda_2+x & \cdots & a+x\\
\vdots & \vdots & \ddots & \vdots\\
b+x & b+x & \cdots & \lambda_n+x
\end{pmatrix}.
\]
\begin{enumerate}[label=(\alph*)]
\item Prove that $p(x)$ is a polynomial of degree at most $1$ with respect to $x$.
\item Calculate the determinant of order $n$ when $x=0$.
\end{enumerate}
\end{problem}

\begin{problem}
For a matrix $A\in M_n(\mathbb{R})$ and two column vectors $u,v\in\mathbb{R}^n$, consider the polynomial $p(x)=\det(A+x\cdot uv^T)$. Prove that $p(x)$ is a polynomial of degree $\le 1$ and find the coefficient of $x$.
\end{problem}

\begin{problem}
Let $A,B\in M_n(\mathbb{R})$. Prove that if $\operatorname{rank}(B)=k$, then the polynomial $p(x)=\det(A+xB)$ has a degree not exceeding $k$.
\end{problem}

\begin{problem}
Let $A$ be a real matrix of order $n$. Given that the sum of all elements in each row equals $\lambda$. Let $J$ be the matrix of all ones. Calculate the determinant $\det(A-tI+J)$ in terms of the characteristic polynomial of $A$, denoted as $P_A(t)=\det(A-tI)$.
\end{problem}

\begin{problem}
Let row matrices $A,B,C\in M_{1\times n}(\mathbb{R})$.
\begin{enumerate}[label=(\alph*)]
\item Prove that if $Bx=0\Rightarrow Ax=0$ for all $x\in\mathbb{R}^n$, then there exists a real number $r$ such that $A=rB$.
\item Prove that if $(Bx=0 \text{ and } Cx=0)\Rightarrow Ax=0$ for all $x\in\mathbb{R}^n$, then there exist real numbers $r,t$ such that $A=rB+tC$.
\end{enumerate}
\end{problem}

\begin{problem}
Let $f,g_1,g_2,\dots,g_k$ be linear functionals on a vector space $V$. Prove that $\displaystyle\bigcap_{i=1}^k \ker(g_i) \subseteq \ker(f)$ if and only if $f\in\operatorname{span}\{g_1,g_2,\dots,g_k\}$.
\end{problem}

\begin{problem}
Let $A\in M_{m\times n}(\mathbb{R})$ and $B\in M_{p\times n}(\mathbb{R})$. Prove that $\ker(A)\subseteq \ker(B)$ if and only if there exists a matrix $P\in M_{p\times m}(\mathbb{R})$ such that $B=P\cdot A$.
\end{problem}

\begin{problem}
Let $f:M_n(\mathbb{R})\to\mathbb{R}$ be a linear functional. Suppose $f(AB)=f(BA)$ for all matrices $A,B\in M_n(\mathbb{R})$. Prove that there exists a real number $c$ such that $f(X)=c\cdot\operatorname{tr}(X)$ for all $X\in M_n(\mathbb{R})$.
\end{problem}

\begin{problem}
Let $A\in M_n(\mathbb{R})$ such that for all vectors $x=(x_1,\dots,x_n)\in\mathbb{R}^n$, the two vectors $xA$ and $x$ are always linearly dependent in $\mathbb{R}^n$. Prove that there exists a real number $a$ such that $A=aI_n$.
\end{problem}

\begin{problem}
Let $A\in M_n(\mathbb{R})$ (with $n\ge 3$). Suppose that for every vector $x\in\mathbb{R}^n$, the set of three vectors $\{x,Ax,A^2x\}$ is always linearly dependent. Prove that the set $\{I,A,A^2\}$ is linearly dependent in the matrix space.
\end{problem}

\begin{problem}
Let $A,B\in M_n(\mathbb{R})$ be matrices such that $\operatorname{rank}(A)\ge 2$. Suppose that for every vector $x\in\mathbb{R}^n$, the vectors $Ax$ and $Bx$ are always linearly dependent. Prove that there exists a real constant $k$ such that $B=kA$.
\end{problem}

\begin{problem}
Let $A=(a_{ij})_{n\times n}$ be a matrix satisfying the strictly diagonally dominant condition: $|a_{ii}| > \sum_{j\ne i}|a_{ij}|$ for all $i=1,\dots,n$. Prove that $A$ is invertible.
\end{problem}

\begin{problem}
Let $A$ be a real skew-symmetric matrix of order $n$.
\begin{enumerate}[label=(\alph*)]
\item If $n$ is odd, prove that $A$ is not invertible.
\item Let $M=I_n+A$. Prove that $M$ is always invertible.
\item Prove that $\det(I+A) \ge 1$ for all real skew-symmetric matrices $A$.
\item Prove that the matrix $Q=(I-A)(I+A)^{-1}$ always has a determinant of $1$ and is an orthogonal matrix.
\end{enumerate}
\end{problem}

\begin{problem}
Let $a_1,\dots,a_n$ and $b_1,\dots,b_n$ be pairwise distinct real numbers such that $a_i+b_j\ne 0$. Prove that the matrix $C$ with elements $c_{ij}=\dfrac{1}{a_i+b_j}$ is invertible.
\end{problem}

\begin{problem}
Let $A$ satisfy $A^k=0$ for some positive integer $k$ (nilpotent matrix). Prove that $I-A$ is always invertible.
\end{problem}

\begin{problem}
Let $A$ be a square matrix of order $n$ satisfying: every main diagonal element $a_{ii}$ is even, and every off-diagonal element $a_{ij}$ ($i\ne j$) is odd. Prove that if $n$ is even, $A$ is invertible over $\mathbb{R}$.
\end{problem}

\begin{problem}
Let $M=\begin{pmatrix}A & B\\ C & D\end{pmatrix}$ be a block matrix where $A$ and $D$ are square matrices. Assume $A$ is invertible. Prove that $M$ is invertible if and only if the Schur complement $S=D-CA^{-1}B$ is invertible.
\end{problem}

\begin{problem}
Let $A$ be an $n\times n$ matrix with elements $a_{ij}=\gcd(i,j)$. Prove that $A$ is invertible.
\end{problem}

\begin{problem}
Let $x_1<x_2<\cdots<x_n$ be distinct real numbers. Consider matrix $A$ with $a_{ij}=|x_i-x_j|$. Prove that for $n>1$, $A$ is always invertible.
\end{problem}

\begin{problem}
Let $C$ be a circulant matrix with first row $(c_0,c_1,\dots,c_{n-1})$. Prove $C$ is invertible if and only if $f(\omega^k)\ne 0$ for all $k=0,\dots,n-1$, where $f(x)=c_0+c_1x+\cdots+c_{n-1}x^{n-1}$ and $\omega=e^{2\pi i/n}$.
\end{problem}

\begin{problem}
Let $A$ and $B$ be two real symmetric positive definite matrices. Prove that their Hadamard product $C=A\circ B$ (where $c_{ij}=a_{ij}b_{ij}$) is also positive definite and thus invertible.
\end{problem}

\begin{problem}
(Sylvester's Rank Inequality) For an $m\times n$ matrix $A$ and an $n\times p$ matrix $B$, prove that
\[
\operatorname{rank}(A)+\operatorname{rank}(B)-n \le \operatorname{rank}(AB) \le \min(\operatorname{rank}(A),\operatorname{rank}(B)).
\]
\end{problem}

\begin{problem}
(Frobenius Rank Inequality) For matrices $A,B,C$ such that the product $ABC$ exists, prove that
\[
\operatorname{rank}(AB)+\operatorname{rank}(BC) \le \operatorname{rank}(B)+\operatorname{rank}(ABC).
\]
\end{problem}

\begin{problem}
Let $P$ be a square matrix satisfying $P^2=P$ (a projection matrix).
\begin{enumerate}[label=(\alph*)]
\item Prove that $\operatorname{rank}(P)+\operatorname{rank}(I-P)=n$.
\item Prove that $\operatorname{rank}(P)=\operatorname{tr}(P)$.
\end{enumerate}
\end{problem}

\begin{problem}
Prove that any $m\times n$ matrix $A$ of rank $r$ can be factored as $A=B\cdot C$, where $B$ is an $m\times r$ matrix (full column rank) and $C$ is an $r\times n$ matrix (full row rank).
\end{problem}

\begin{problem}
Let $A$ be a square matrix and $u,v$ be column vectors. Prove that $\operatorname{rank}(A+uv^T)$ can only take the values: $\operatorname{rank}(A)-1$, $\operatorname{rank}(A)$, or $\operatorname{rank}(A)+1$.
\end{problem}

\begin{problem}
Let $A,B$ be $n\times n$ matrices. Prove that $\operatorname{rank}(A+B)=\operatorname{rank}(A)+\operatorname{rank}(B)$ if and only if $\operatorname{Im}(A)\cap \operatorname{Im}(B)=\{0\}$ and $\operatorname{Im}(A^T)\cap \operatorname{Im}(B^T)=\{0\}$.
\end{problem}

\begin{problem}
Let $A\in\mathbb{R}^{m\times n}$ and $B\in\mathbb{R}^{n\times p}$. Prove that $\operatorname{rank}(AB)=\operatorname{rank}(B)-\dim(\operatorname{Im}(B)\cap \ker(A))$. Deduce the necessary and sufficient condition for $\operatorname{rank}(AB)=\operatorname{rank}(B)$.
\end{problem}

\begin{problem}
Let $A$ be a real matrix satisfying $A^3=A$. Prove that $\operatorname{rank}(A)=\operatorname{tr}(A^2)$.
\end{problem}

\begin{problem}
Calculate the determinant of order $n$:
\[
D = \det\begin{pmatrix}
1+x_1 & 1 & 1 & \cdots & 1\\
1 & 1+x_2 & 1 & \cdots & 1\\
1 & 1 & 1+x_3 & \cdots & 1\\
\vdots & \vdots & \vdots & \ddots & \vdots\\
1 & 1 & 1 & \cdots & 1+x_n
\end{pmatrix}.
\]
\end{problem}

\begin{problem}
Calculate the determinant of order $n$ (n $\ge 3$) with elements $a_{ij}=\cos(x_i+y_j)$.
\end{problem}

\begin{problem}
Let $S=\{x_1,\dots,x_n\}$ be a set of $n$ distinct positive integers, closed under the GCD operation. Prove that $\det(\gcd(x_i,x_j)) = \prod_{i=1}^n \varphi(x_i)$.
\end{problem}

\begin{problem}
Let $A$ and $B$ be square matrices of order $n$. If $\operatorname{rank}(B)=1$, prove that the function $f(x)=\det(A+xB)$ is a polynomial of at most degree $1$ in $x$.
\end{problem}

\begin{problem}
Let $A$ be an $n\times n$ matrix with $A_{ii}=\alpha$ and $A_{ij}=\beta$ for $i\ne j$. Calculate $\det(A)$.
\end{problem}

\begin{problem}
Let $f:V\to V$ be an endomorphism on a finite-dimensional vector space $V$.
\begin{enumerate}[label=(\alph*)]
\item Prove there exists a positive integer $k$ such that $\ker(f^k)=\ker(f^{k+1})$ and $\operatorname{Im}(f^k)=\operatorname{Im}(f^{k+1})$.
\item With such $k$, prove that $V=\ker(f^k)\oplus \operatorname{Im}(f^k)$.
\end{enumerate}
\end{problem}

\begin{problem}
Let $f:V\to W$ and $g:V\to U$ be linear mappings. Prove there exists a linear mapping $h:W\to U$ such that $g=h\circ f$ if and only if $\ker(f)\subseteq \ker(g)$.
\end{problem}

\begin{problem}
Let $f:V\to V$ satisfy $f^2=0$. Prove that $\operatorname{Im}(f)\subseteq \ker(f)$. Furthermore, for an arbitrary $f$, prove that $\dim(\operatorname{Im}(f)\cap \ker(f)) = \dim(\operatorname{Im}(f))-\dim(\operatorname{Im}(f^2))$.
\end{problem}

\begin{problem}
Let $V$ be an $n$-dimensional vector space. Let $H_1,H_2,\dots,H_k$ be $k$ hyperplanes (subspaces of dimension $n-1$). Prove that $\dim(H_1\cap H_2\cap\cdots\cap H_k) \ge n-k$.
\end{problem}

\begin{problem}
Let $f:U\to V$ and $g:V\to W$ be linear mappings. Prove the inequality: $\dim(\ker(g\circ f)) \le \dim(\ker(f))+\dim(\ker(g))$.
\end{problem}

\begin{problem}
Let $A,B$ be square matrices of order $n$ over $\mathbb{C}$.
\begin{enumerate}[label=(\alph*)]
\item Prove that the equality $AB-BA=I$ is impossible.
\item Suppose $AB-BA=A$. Prove that $A$ is nilpotent.
\end{enumerate}
\end{problem}

\begin{problem}
Let $P,Q$ be two projections ($P^2=P$, $Q^2=Q$). Prove that $\operatorname{tr}(P-Q)^3=\operatorname{tr}(P-Q)$.
\end{problem}

\begin{problem}
Let $f:V\to V$ be an endomorphism. Prove that $\operatorname{tr}(f)=0$ if and only if there exist two endomorphisms $g,h:V\to V$ such that $f=g\circ h - h\circ g$.
\end{problem}

\begin{problem}
Let $A$ and $B$ be diagonalizable matrices of order $n$. Prove that $A$ and $B$ commute ($AB=BA$) if and only if they are simultaneously diagonalizable.
\end{problem}

\begin{problem}
Let $A$ be a matrix over $\mathbb{C}$ with $n$ distinct eigenvalues. Prove that a matrix $B$ commutes with $A$ if and only if $B$ is a polynomial in $A$.
\end{problem}

\begin{problem}
Let $A$ be an $n\times n$ matrix over $\mathbb{C}$. Consider the linear mapping $\phi_A(X)=AX-XA$. Prove that $\phi_A$ is diagonalizable if and only if $A$ is diagonalizable.
\end{problem}

\begin{problem}
Let $A\in M_n(\mathbb{Q})$ satisfy $A^p=I$ where $p$ is prime. Prove that $A$ is diagonalizable over $\mathbb{Q}$ if and only if $A=I$.
\end{problem}

\begin{problem}
Prove that a complex matrix $A$ is diagonalizable if and only if for all $\lambda\in\mathbb{C}$, $\operatorname{rank}(A-\lambda I) = \operatorname{rank}((A-\lambda I)^2)$.
\end{problem}

\begin{problem}
Let $A$ be a complex normal matrix ($AA^*=A^*A$). Prove it is unitarily diagonalizable.
\end{problem}

\begin{problem}
Let $A$ be a complex matrix of order $n$. Prove that $A$ is nilpotent ($A^n=0$) if and only if $\operatorname{tr}(A^k)=0$ for all $k=1,2,\dots,n$.
\end{problem}

\begin{problem}
The Contrast Matrix of order $n\ge 2$, denoted $K_n$, is defined by $K_{ij}=\dfrac{i}{j}-\dfrac{j}{i}$.
\begin{enumerate}[label=(\alph*)]
\item Prove $K_n$ is skew-symmetric with trace $0$.
\item Prove $\operatorname{rank}(K_n)=2$.
\item Prove the characteristic polynomial has the form $P_K(\lambda)=\lambda^{n-2}(\lambda^2+\Delta_n)$.
\end{enumerate}
\end{problem}

\begin{problem}
Let $S(n,m)$ be the set of matrices where the sum of even-indexed columns equals the sum of odd-indexed columns. If $A\in S(n,m)$, find the smallest eigenvalue of the matrix $M=I+A^TA$.
\end{problem}

\begin{problem}
Let $A$ be an $n\times m$ real matrix with $n<m$. Prove that $\lambda=1$ is an eigenvalue of $M=I_m+A^TA$ with algebraic multiplicity at least $m-n$.
\end{problem}

\begin{problem}
Let $A\in M_{2\times 2}(\mathbb{C})$ have two distinct eigenvalues. Let
 $S=\left\{N\in M_{2\times 2}(\mathbb{C}) \;\middle|\; \begin{pmatrix}A & N\\ 0 & A\end{pmatrix} \text{ is diagonalizable}\right\}$.
Prove that $S$ is a $2$-dimensional subspace of $M_{2\times 2}(\mathbb{C})$.
\end{problem}

\begin{problem}
Let $a_1,\dots,a_n\in\mathbb{C}\setminus\{0\}$ and $A=\left(\dfrac{a_i}{a_j}\right)_{1\le i,j\le n}$. Prove $A$ is diagonalizable and find an invertible matrix $C$ such that $C^{-1}AC$ is diagonal.
\end{problem}

\begin{problem}
Let $A=(a_{ij})\in M_{n\times n}(\mathbb{R})$ satisfy $a_{ii}=0$ for all $i$ and $a_{ij}>0$ for all $i\ne j$. Find the minimum possible value of $\operatorname{rank}(A)$.
\end{problem}

\begin{problem}
Let $V$ be a $k$-dimensional real vector space with standard inner product. Given a system of $n$ vectors $x_1,x_2,\dots,x_n\in V$, consider the matrix $D\in M_n(\mathbb{R})$ defined by the squared distances $D_{ij}=\|x_i-x_j\|^2$. Prove that $\operatorname{rank}(D)\le k+2$.
\end{problem}

\begin{problem}
Calculate the following determinant of order $n$:
\[
A_n=\begin{pmatrix}
2021 & 1 & 1 & \cdots & 1 & 1\\
1 & 2021 & 1 & \cdots & 1 & 1\\
1 & 1 & 2021 & \cdots & 1 & 1\\
\vdots & \vdots & \vdots & \ddots & \vdots & \vdots\\
1 & 1 & 1 & \cdots & 2021 & 1\\
1 & 1 & 1 & \cdots & 1 & 2021
\end{pmatrix}.
\]
\end{problem}

\begin{problem}
Calculate the following determinant of order $n$:
\[
B_n=\begin{pmatrix}
5 & 3 & 0 & 0 & \cdots & 0 & 0 & 0\\
2 & 5 & 3 & 0 & \cdots & 0 & 0 & 0\\
0 & 2 & 5 & 3 & \cdots & 0 & 0 & 0\\
\vdots & \vdots & \vdots & \vdots & \ddots & \vdots & \vdots & \vdots\\
0 & 0 & 0 & 0 & \cdots & 5 & 3 & 0\\
0 & 0 & 0 & 0 & \cdots & 2 & 5 & 3\\
0 & 0 & 0 & 0 & \cdots & 0 & 2 & 5
\end{pmatrix}.
\]
\end{problem}

\begin{problem}
Calculate the following determinant of order $n$, where $a,b\in\mathbb{R}$:
\[
C_n=\begin{pmatrix}
a+b & ab & 0 & 0 & \cdots & 0 & 0 & 0\\
1 & a+b & ab & 0 & \cdots & 0 & 0 & 0\\
0 & 1 & a+b & ab & \cdots & 0 & 0 & 0\\
\vdots & \vdots & \vdots & \vdots & \ddots & \vdots & \vdots & \vdots\\
0 & 0 & 0 & 0 & \cdots & a+b & ab & 0\\
0 & 0 & 0 & 0 & \cdots & 1 & a+b & ab\\
0 & 0 & 0 & 0 & \cdots & 0 & 1 & a+b
\end{pmatrix}.
\]
\end{problem}

\begin{problem}
Let $M$ be an $n\times n$ matrix where every diagonal entry is $a$ and every off-diagonal entry is $b$:
\[
M=\begin{pmatrix}
a & b & b & \cdots & b\\
b & a & b & \cdots & b\\
b & b & a & \cdots & b\\
\vdots & \vdots & \vdots & \ddots & \vdots\\
b & b & b & \cdots & a
\end{pmatrix}.
\]
Show that
\[
M^{-1} = \frac{1}{a-b}I - \frac{b}{(a-b)(a+(n-1)b)}J
\]
\end{problem}

\begin{problem}
Let $K$ be a field of characteristic zero and $A\in M_{n\times n}(K)$ be the permutation
matrix associated to a permutation $\sigma\in S_n$ (here $n$ is a positive integer).
Suppose the cycle type of $\sigma$ is $(k_1,k_2,\ldots,k_n)$, where $k_i$ is the number
of $i$-cycles in $\sigma$ for $1\le i\le n$.

Now consider the linear transformations
\[
L_A:M_{n\times n}(K)\longrightarrow M_{n\times n}(K),
\qquad
R_A:M_{n\times n}(K)\longrightarrow M_{n\times n}(K)
\]
given by $L_A(B)=AB$, the left multiplication of $B$ by $A$, and
$R_A(B)=BA$, the right multiplication of $B$ by $A$.
Note here $L_A,R_A$ are linear transformations on the $n^2$-dimensional vector space.
Find $\operatorname{Rank}(L_A+R_A)$.
\end{problem}

\begin{problem}
Let $K$ be a field of characteristic zero and $A\in M_{n\times n}(K)$ be any matrix.
Consider the linear transformations
\[
L_A:M_{n\times n}(K)\longrightarrow M_{n\times n}(K),
\qquad
R_A:M_{n\times n}(K)\longrightarrow M_{n\times n}(K)
\]
given by $L_A(B)=AB$ and $R_A(B)=BA$.

Show that the space
\[
S=\ker(L_A+R_A)\cap\operatorname{Image}(L_A+R_A)
\]
consists of nilpotent matrices; that is, any matrix $B\in S$ satisfies that
\[
B^n=0.
\]
\end{problem}

\begin{problem}
Let $F$ be a field and let
\[
f(X)=\prod_{i=1}^{k}(X-\lambda_i)^{r_i},\qquad
\widetilde f(X)=\prod_{i=1}^{k}(X-\lambda_i)^{\widetilde r_i},
\]
and
\[
g(X)=\prod_{i=1}^{k}(X-\lambda_i)^{s_i},\qquad
\widetilde g(X)=\prod_{i=1}^{k}(X-\lambda_i)^{\widetilde s_i},
\]
where $\lambda_1,\lambda_2,\ldots,\lambda_k$ are distinct elements in $F$,
$r_i\le \widetilde r_i$ and $s_i\le \widetilde s_i$, $1\le i\le k$,
are nonnegative integers.

Let $V$ be a vector space of dimension $n+m$ and $W\subseteq V$ be a subspace
of dimension $n$ where
\[
n=\widetilde r_1+\cdots+\widetilde r_k,\qquad
m=\widetilde s_1+\cdots+\widetilde s_k.
\]
Let $T:V\to V$ be a linear transformation with $T(W)\subseteq W$ and suppose the
minimal polynomial of $T|_W$ is $f(X)$, the characteristic polynomial of $T|_W$
is $\widetilde f(X)$, and the minimal polynomial of the induced linear transformation
\[
\widetilde T:V/W\longrightarrow V/W
\]
is $g(X)$, while its characteristic polynomial is $\widetilde g(X)$.

Let $t_1,t_2,\ldots,t_k$ be nonnegative integers. Show that there exists a linear
transformation $T$ on $V$ and an invariant subspace $W$ as above with $T$ on $V$
having minimal polynomial
\[
h(X)=\prod_{i=1}^{k}(X-\lambda_i)^{t_i}
\]
if and only if the given nonnegative integers $t_1,t_2,\ldots,t_k$ satisfy
\[
\max(r_i,s_i)\le t_i\le r_i+s_i,\qquad 1\le i\le k.
\]
\end{problem}

\begin{problem}
For an integer $n\ge2$, let $A$ and $B$ be singular $n\times n$ complex matrices such that
\[
(AB)^n=0.
\]
Prove that at least one of $(AB)^{n-1}$ and $(BA)^{n-1}$ is the zero matrix.
\end{problem}

\begin{problem}
Let $A$ and $B$ be $n\times n$ complex matrices such that
\[
\operatorname{rank}(AB)=\operatorname{rank}(BA).
\]
Prove that
\[
AB^2A=AB
\]
if and only if
\[
BA^2B=BA.
\]
\end{problem}

\begin{problem}
Let $S$ be a finite set of real numbers, and let $T$ be the set of all $n\times n$
matrices having entries in $S$. Prove
\[
\sum_{A\in T}\operatorname{trace}(A^2)
=
\sum_{A\in T}(\operatorname{trace}(A))^2.
\]
\end{problem}

\begin{problem}
For any $n\times n$ complex matrix $M$, prove
\[
\operatorname{rank}(M)+\operatorname{rank}(M-M^3)
=
\operatorname{rank}(M-M^2)+\operatorname{rank}(M+M^2).
\]
\end{problem}

\begin{problem}
Let $M_n$ be the set of $n\times n$ complex matrices. We write $I$ for the identity
matrix in $M_n$. Suppose that $A,B\in M_n$ satisfy
\[
(I-A)^2B+2(I-A)B(I-A)+B(I-A)^2=B,
\]
and suppose that there are nonzero complex numbers $\alpha,\beta$ and odd positive
integers $j,k$ such that
\[
A^j=\alpha I,\qquad B^k=\beta I.
\]
Prove that the following two conditions are equivalent:
\begin{enumerate}
    \item There exists a finite subset $S$ of $M_n$ containing $B$ such that
    \[
    AX+XA\in S
    \]
    for all $X\in S$.
    \item The matrix $A$ is $\frac12 I$.
\end{enumerate}
\end{problem}

\begin{problem}
A derangement is a permutation with no fixed points. Let $D_n$ be the set of derangements
on $\{1,2,\ldots,n\}$. For $\sigma\in D_n$, let
\[
\omega(\sigma)=\frac1n\prod_{i=1}^{n}(\sigma(i)-i)
\]
and
\[
\lambda_n=(-1)^{\lfloor n/2\rfloor}\sum_{\sigma\in D_n}\omega(\sigma).
\]
Prove that $\lambda_n=0$ when $n$ is odd and
\[
\lambda_n=\frac{q_n}{2^n},
\]
where $q_n$ is rational and $q_n=1$ when $n=2$ and $q_n>1$ when $n>2$ and is even.
\end{problem}

\begin{problem}
Let $k$ be an integer with $k\ge2$. In terms of $k$, find the smallest positive integer
$n$ such that there exists an $n\times n$ integer matrix $A$ satisfying the following
two conditions:
\begin{enumerate}
    \item $A^TA=AA^T$;
    \item $k$ is the least positive integer $m$ such that $A^m$ equals the identity matrix.
\end{enumerate}
\end{problem}

\begin{problem}
Let $A$ and $B$ be complex $n\times n$ matrices, and suppose
\[
A^mB=A
\]
for some integer $m$ with $m\ge2$. Prove
\[
A^mB
=
A^{m-1}BA
=
A^{m-2}BA^2
=
\cdots
=
ABA^{m-1}.
\]
\end{problem}

\begin{problem}
A matrix $A$ with $i,j$-entry $a_{ij}$ is special if
\[
a_{ij}a_{sk}-a_{is}a_{jk}+a_{ik}a_{js}=0
\]
for all $i,j,s,k$ such that
\[
1\le i<j<s<k\le n.
\]
A matrix $A$ is strong if all its $3\times3$ submatrices have rank $3$.

\begin{enumerate}
    \item Prove that, when $n\ge6$, no $n\times n$ real matrix is special and strong.
    \item Show that there is a $5\times5$ matrix that is special and strong.
\end{enumerate}
\end{problem}

\begin{problem}
Let $A$ and $B$ be $n\times n$ matrices that are real, symmetric, and positive definite.
Find all $n\times n$ matrices $X$ that are real, symmetric, and positive definite and that
satisfy
\[
XAX=BX^{-1}B.
\]
\end{problem}

\begin{problem}
Let $n$ and $k$ be integers with $1\le k\le n$, and let $A$ be an $n\times n$ matrix
with $i,j$-entry given by
\[
a_{ij}=
\begin{cases}
1, & \text{if } j\equiv i+1\pmod{k},\\
0, & \text{otherwise}.
\end{cases}
\]
What is the characteristic polynomial and minimum polynomial of $A$?
\end{problem}

\begin{problem}
Let $A^*$ denote the conjugate transpose of a complex matrix $A$. Show that
\[
A^5=AA^*AA^*A
\]
if and only if
\[
A=B+C
\]
for some matrices $B$ and $C$ satisfying
\[
B=B^*,\qquad C=-C^*,\qquad BC=CB=0.
\]
\end{problem}

\begin{problem}
Find all prime numbers $p$ such that there exists a $2\times2$ matrix $A$ with integer
entries, other than the identity matrix $I$, for which
\[
A^p+A^{p-1}+\cdots+A=pI.
\]
\end{problem}

\begin{problem}
Let $k$ and $n$ be positive integers with $k<n$. Characterize the $n\times n$ real matrices
$M$ with the property that for all $v\in\mathbb R^n$ with at most $k$ nonzero entries,
$Mv$ also has at most $k$ nonzero entries.
\end{problem}

\begin{problem}
Let $A$ be an $n\times n$ positive-definite Hermitian matrix, with minimum and maximum
eigenvalues $\lambda$ and $\omega$, respectively. Prove
\[
\left(
\frac{1}{\omega}\frac{\operatorname{Tr}(A)}{n}
+\frac{\omega n}{\operatorname{Tr}(A)}
\right)^n
\le
\det\left(\frac{1}{\omega}A+\omega A^{-1}\right),
\]
and
\[
\left(
\frac{1}{\lambda\,\operatorname{Tr}(A^{-1})}
+\frac{\lambda\,\operatorname{Tr}(A^{-1})}{n}
\right)^n
\le
\det\left(\frac{1}{\lambda}A+\lambda A^{-1}\right).
\]
\end{problem}

\begin{problem}
Let $A$ be an $n\times n$ skew-symmetric real matrix. Show that for positive real numbers
$x_1,\ldots,x_k$ with $k\ge2$,
\[
\det(A+x_1I)\cdots\det(A+x_kI)
\ge
\left(
\det\left(A+(x_1\cdots x_k)^{1/k}I\right)
\right)^k.
\]
In addition, show that if also all $x_i$ lie on the same side of $1$, then
\[
\det(A+I)^{k-1}
\det\left(A+x_1\cdots x_k I\right)
\ge
\det(A+x_1I)\cdots\det(A+x_kI).
\]
\end{problem}

\begin{problem}
\textbf{Definitions.}
\begin{enumerate}[label=\alph*)]
\item A matrix is called a \emph{complex matrix} if all its entries are complex numbers. (In particular, any vector in $\mathbb{C}^n$ is considered an $n\times 1$ complex matrix.)
\item For any complex matrix $A$, we denote the matrix $\overline{A}^T$ by $A^*$.
\item Let $U_2(\mathbb{C})$ be the group of all unitary $2\times 2$ complex matrices. In other words,
\[
U_2(\mathbb{C}) = \{U \in GL_2(\mathbb{C}) \mid U^* = U^{-1}\}.
\]
\item A complex matrix $A$ is called \emph{Hermitian} if it satisfies $A^*=A$.
\end{enumerate}
\textbf{Problem.} Let $A_1, A_2, \dots, A_n$ be $n$ Hermitian $2\times 2$ complex matrices. Define a function $F$ from the Cartesian product $(U_2(\mathbb{C}))^n$ to $\mathbb{R}$ by
\[
F(U_1,U_2,\dots,U_n) = \det\!\left(\sum_{k=1}^n U_k^* A_k U_k\right) \qquad \text{for every } (U_1,\dots,U_n)\in (U_2(\mathbb{C}))^n.
\]
Show that
\[
\min_{U\in (U_2(\mathbb{C}))^n} F(U) = \sum_{k=1}^n \sigma_1(A_k)\cdot \sum_{k=1}^n \sigma_2(A_k),
\]
where $\sigma_1(A_j)$ and $\sigma_2(A_j)$ denote the greatest and the least eigenvalue of the matrix $A_j$, respectively, for every $j\in\{1,2,\dots,n\}$.
\end{problem}

\begin{problem}
\begin{enumerate}[label=(\alph*)]
\item Show that for any $m\times n$ complex matrix $E=(e_{ij})$,
\[
\det(E^*E) \le \prod_{j=1}^n \sum_{i=1}^m |e_{ij}|^2.
\]
\item Let $F$ be an $m\times n$ matrix. If $G$ consists of some columns of $F$ and $H$ contains the remaining columns of $F$, show that
\[
\det(F^*F) \le \det(G^*G)\det(H^*H).
\]
\end{enumerate}
\end{problem}

\begin{problem}
Let $X$ and $Y$ be square matrices of the same size. Show that
\[
|\det(I-YX^*)|^2 \le \det(I+XX^*)\det(I+Y^*Y)
\]
and
\[
|\det(X+Y)|^2 \le \det(I+XX^*)\det(I+Y^*Y).
\]
\end{problem}

\begin{problem}
Let $A\in M_n(\mathbb{C})$. Show that
\begin{enumerate}[label=(\alph*)]
\item $A=AA^*S$ for some matrix $S$.
\item If $A$ is nonsingular, then matrix $S$ in (a) is unique.
\item $A=(AA^*)^{1/2}T$ for some matrix $T$.
\item If $A$ is nonsingular, then $T$ in (c) is unique and it is unitary.
\end{enumerate}
\end{problem}

\begin{problem}
Let $\{\beta_1,\dots,\beta_n\}$ be a basis of a vector space $V$ of dimension $n$ and let $\{\alpha_1,\dots,\alpha_m\}$ be a set of vectors in $V$. Write each $\alpha_i$ as
\[
\alpha_i = \sum_{j=1}^n a_{ij}\beta_j, \qquad i=1,\dots,m.
\]
Let $A=(a_{ij})$ be the $m\times n$ coefficient matrix and $\{i_1,\dots,i_k\}\subseteq\{1,\dots,m\}$. Show that $\{\alpha_{i_1},\dots,\alpha_{i_k}\}$ is linearly independent if and only if the corresponding rows $i_1,\dots,i_k$ of $A$ are linearly independent, and if and only if there is a nonsingular $k\times k$ submatrix of $A$ lying in rows $i_1,\dots,i_k$ of $A$. Conclude that $\dim \operatorname{Span}\{\alpha_1,\dots,\alpha_m\} = r(A)$.
\end{problem}

\begin{problem}
Let $\alpha=\{\alpha_1,\dots,\alpha_m\}$ and $\beta=\{\beta_1,\dots,\beta_n\}$ be two sets of vectors in a vector space $V$ such that each $\alpha_i$ is a linear combination of $\beta_1,\dots,\beta_n$:
\[
\alpha_i = \sum_{j=1}^n a_{ij}\beta_j, \qquad i=1,\dots,m.
\]
Let $A=(a_{ij})$ be the $m\times n$ coefficient matrix. Show that
\begin{enumerate}[label=(\alph*)]
\item $\dim(\operatorname{Span}\alpha) \le \dim(\operatorname{Span}\beta)$.
\item $\dim(\operatorname{Span}\alpha) \le r(A)$.
\item If $\beta$ is linearly independent, then $\dim(\operatorname{Span}\alpha)=r(A)$.
\item If $\alpha$ is linearly independent, then $r(A)=m$ and $n\ge m$.
\end{enumerate}
\end{problem}

\begin{problem}
Let $A$ be an $n\times n$ real matrix.
\begin{enumerate}[label=(\alph*)]
\item If $A^t=-A$ and $n$ is odd, show that $|A|=0$.
\item If $A^2+I=0$, show that $n$ must be even.
\item Do (a) and (b) remain true for complex matrices?
\end{enumerate}
\end{problem}

\begin{problem}
Let $A\in M_n(\mathbb{C})$. If $AA^t=I$ and $|A|<0$, show that $|A+I|=0$.
\end{problem}

\begin{problem}
If $A,B,C$, and $D$ are $n\times n$ matrices such that $ABCD=I$, show that
\[
ABCD = DABC = CDAB = BCDA = I.
\]
What does $D^{-1}C^{-1}B^{-1}$ equal? Is it true that $BCAD=I$?
\end{problem}

\begin{problem}
Let $n\in\mathbb{N}$, $n\ge 2$.
\begin{enumerate}[label=(\alph*)]
\item Prove that $\det(A^2-B^2)(C^2-B^2)\ge 0$ for all $A,B,C\in M_n(\mathbb{R})$ with $AB=BC$.
\item Find all values $k\ge 1$ such that $\det(A^k-B^2)(C^k-B^2)\ge 0$ holds for all $A,B,C\in M_n(\mathbb{R})$ with $AB=BC$.
\end{enumerate}
\end{problem}

\begin{problem}
Let $n\in\mathbb{N}$, $n\ge 2$. Find all the values $k\in\mathbb{N}$, $k\ge 1$, for which the following statement holds:
\[
\text{``If } A\in M_n(\mathbb{C}) \text{ is such that } A^kA^*=A, \text{ then } A=A^*.\text{''} \tag{$*$}
\]
(Here, $A^*=\overline{A}^t$ denotes the transpose conjugate of $A$.)
\end{problem}

\begin{problem}
Let $\operatorname{adj}(A)$ denote the adjugate matrix (or cofactor matrix) of $A=(a_{ij})\in M_n(\mathbb{C})$, that is, $\operatorname{adj}(A)$ is the $n\times n$ matrix whose $(i,j)$-entry is the cofactor $(-1)^{i+j}|A_{ji}|$ of $a_{ji}$, where $A_{ji}$ is the submatrix obtained from $A$ by deleting the $j$-th row and the $i$-th column. Show that
\begin{enumerate}[label=(\alph*)]
\item $r(A)=n$ if and only if $r(\operatorname{adj}(A))=n$.
\item $r(A)=n-1$ if and only if $r(\operatorname{adj}(A))=1$.
\item $r(A)<n-1$ if and only if $r(\operatorname{adj}(A))=0$.
\item $|\operatorname{adj}(A)| = |A|^{n-1}$.
\item $\operatorname{adj}(\operatorname{adj}(A)) = |A|^{n-2}A$.
\item $\operatorname{adj}(AB) = \operatorname{adj}(B)\operatorname{adj}(A)$.
\item $\operatorname{adj}(XAX^{-1}) = X(\operatorname{adj}(A))X^{-1}$ for any invertible $X\in M_n(\mathbb{C})$.
\item $|\underbrace{\operatorname{adj}\cdots\operatorname{adj}}_{k}(A)| = |A|$ when $A$ is $2\times 2$.
\item If $A$ is Hermitian, so is $\operatorname{adj}(A)$.
\item Find a formula for $\underbrace{\operatorname{adj}\cdots\operatorname{adj}}_{k}(A)$ when $|A|=1$.
\end{enumerate}
\end{problem}

\begin{problem}
Let $A\in M_n(\mathbb{C})$. Show that
\begin{enumerate}[label=(\alph*)]
\item $\operatorname{Im}(AA^*)=\operatorname{Im}A$. Does $\operatorname{Im}(A^*A)=\operatorname{Im}A$?
\item $A=AA^*B$ for some matrix $B\in M_n(\mathbb{C})$.
\item If $A^*A=AA^*$ and $A^2x=0$, then $Ax=0$.
\end{enumerate}
\end{problem}

\begin{problem}
Show that $A$ is invertible if $A=(a_{ij})\in M_n(\mathbb{C})$ satisfies, for each $i$,
\[
|a_{ii}| > \sum_{j=1,\, j\ne i}^n |a_{ij}|. \qquad \text{(Row diagonal dominance)}
\]
Show that $I-B$ is invertible if $B=(b_{ij})\in M_n(\mathbb{C})$ satisfies
\[
1 > \sum_{j=1}^n |b_{ij}|, \qquad i=1,2,\dots,n.
\]
\end{problem}

\begin{problem}
Find all matrices $A \in M_n(\mathbb{Z})$ for which
\[
2025A^{2025} = A^{2024} + A^{2023} + \cdots + A.
\]
\end{problem}

\begin{problem}
Let
\[
x_{ij} =
\begin{cases}
a_i & \text{if } i = j, \\
0 & \text{if } i \neq j,\ i+j \neq 2n+1, \\
b_i & \text{if } i+j = 2n+1,
\end{cases}
\]
where $a_i, b_i$ are real numbers. Evaluate the determinant $\Delta_{2n} = |x_{ij}|$.
\end{problem}

\begin{problem}
Let $n \geq 2$ be an integer, and let $A$ be an $n \times n$ real matrix with minimal polynomial $x^n + x - 1$. Show that
\[
\operatorname{tr}\!\left((nA^{n-1}+I)^{-1}A^{n-2}\right) = 0.
\]
(Here $\operatorname{tr}(B) = \sum_{i=1}^n b_{ii}$ for $B = (b_{ij})_{i,j=1}^n$.)
\end{problem}

\begin{problem}
Let $A, B$ be real square matrices satisfying
\[
A^{2001} = 0, \qquad AB = A + B.
\]
Prove that $\det B = 0$.
\end{problem}

\begin{problem}
Denote by $(a,b)$ the inner product of two vectors $a, b \in \mathbb{R}^n$. Let $a_1, a_2, \dots, a_k \in \mathbb{R}^n$, and put
\[
A =
\begin{bmatrix}
(a_1,a_1) & (a_1,a_2) & \cdots & (a_1,a_{k-1}) & (a_1,a_k) \\
(a_2,a_1) & (a_2,a_2) & \cdots & (a_2,a_{k-1}) & (a_2,a_k) \\
\vdots & \vdots & \ddots & \vdots & \vdots \\
(a_{k-1},a_1) & (a_{k-1},a_2) & \cdots & (a_{k-1},a_{k-1}) & (a_{k-1},a_k) \\
(a_k,a_1) & (a_k,a_2) & \cdots & (a_k,a_{k-1}) & (a_k,a_k)
\end{bmatrix}.
\]
Prove that
\begin{enumerate}[label=\arabic*)]
\item $\det A \geq 0$;
\item $A$ is a symmetric matrix and all of its eigenvalues are nonnegative.
\end{enumerate}
\end{problem}

\begin{problem}
Let $A$ be a square matrix of order $n$ all of whose entries are even integers. Prove that $A$ cannot have an odd integer as an eigenvalue.
\end{problem}

\begin{problem}
Let $A,B\in M_n(\mathbb{R})$ be matrices. Consider the matrix function $f:M_n(\mathbb{C})\to M_n(\mathbb{C})$ defined by
\[
f(Z)=AZ+B\bar{Z},\quad Z\in M_n(\mathbb{C}),
\]
where $\bar{Z}$ is the matrix having as elements the conjugates of the elements of $Z$. Prove that the following statements are equivalent:
\begin{enumerate}[label=(\roman*)]
\item $f$ is injective;
\item $f$ is surjective;
\item the matrices $A+B$ and $A-B$ are invertible.
\end{enumerate}
\end{problem}

\begin{problem}
Fix $n\in\mathbb{N}$. Let $a_1,a_2,\ldots,a_n$ be some real numbers. Define $A\in M_n(\mathbb{R})$ by $A=[a_{ij}]_{1\le i,j\le n}$ where
\[
a_{ij}=\delta_{ij}+a_i+a_j,\quad 1\le i,j\le n.
\]
($\delta_{ij}=1$ if $i=j$ and $\delta_{ij}=0$ if $i\ne j$.) Compute the characteristic polynomial of $A$. [Caution: The value of $n$ matters in these computations.]
\end{problem}

\begin{problem}
Let $A \in \mathcal{M}_n(\mathbb{R})$ be a symmetric matrix and let $k$ be a positive even integer such that
\[ [\mathrm{trace}(A^k)]^{k+1} = [\mathrm{trace}(A^{k+1})]^k. \]
Prove that $A^n = \mathrm{trace}(A) A^{n-1}$. If $k$ is a positive odd integer, does the conclusion still hold?
\end{problem}

\begin{problem}
Let $A, B \in \mathcal{M}_n(\mathbb{R})$ be commuting matrices. Let $\alpha \in \mathbb{C}$ be an eigenvalue of $A+B$. Prove that there is a decomposition $\alpha = \lambda + \mu$ where $\lambda \in \mathbb{C}$ is an eigenvalue of $A$ and $\mu \in \mathbb{C}$ is an eigenvalue of $B$.
\end{problem}

\begin{problem}
Let $A \in \mathcal{M}_n(\mathbb{R})$ with $\mathrm{rank}(A) = r \geq 1$. Prove that $A^2 = A$ if and only if there exist matrices $B \in \mathcal{M}_{n\times r}(\mathbb{R})$ and $C \in \mathcal{M}_{r\times n}(\mathbb{R})$, both of rank $r$, such that $A = BC$ and $CB = I_r$. Moreover, show that if $A^2=A$ then $\det(I_n+A) = 2^r$ and $\det(2I_n - A) = 2^{n-r}$.
\end{problem}

\begin{problem}
Let $A,B \in \mathcal{M}_n(\mathbb{R})$ be invertible matrices satisfying $AB+BA = O_n$. Prove that the matrices $I_n, A, B, AB$ are linearly independent.
\end{problem}

\begin{problem}
Let $A,B \in \mathcal{M}_n(\mathbb{R})$ satisfy $A^{17}=B^{16}=I_n$ and $AB=BA$. Prove that $I_n+A+B$ is invertible.
\end{problem}

\begin{problem}
Let $A=[a_{ij}] \in \mathcal{M}_n(\mathbb{R})$, with $n$ odd, satisfy:
\begin{enumerate}[label=\alph*)]
\item $a_{ii}=\lambda$,
\item $a_{ij}=-a_{ji}$ for all $i\neq j$.
\end{enumerate}
Find the condition on $\lambda$ so that the system $Ax=b$ has a unique solution, where $x=[x_1,\dots,x_n]^t$, $b=[b_1,\dots,b_n]^t$.
\end{problem}

\begin{problem}
Let $A,B \in \mathcal{M}_3(\mathbb{Z})$ satisfy
\[ AB = \begin{bmatrix} 1 & 2k & k(2k+1) \\ 0 & 1 & 2k \\ 0 & 0 & 1 \end{bmatrix},\quad k \in \mathbb{N}. \]
Prove that there exists $C \in \mathcal{M}_3(\mathbb{Z})$ such that $BA = C^k$.
\end{problem}

\begin{problem}
Let $A=[a_{ij}]$ be an $n\times n$ matrix satisfying $I+A+A^2+\cdots+A^{2016}=0$. Prove that $A^2-I$ is invertible.
\end{problem}

\begin{problem}
Let $M = \begin{pmatrix}1 & -1\\ 1& -1\end{pmatrix}$.
\begin{enumerate}[label=\alph*)]
\item Prove that if $A$ is a real $2\times2$ matrix with $AM=MA$, then there exist real numbers $\alpha,\beta$ such that $A=\alpha I+\beta M$, where $I$ is the $2\times2$ identity matrix.
\item Using part a), find all real $2\times 2$ matrices $X$ such that $X^{2016}-X^{2010}=6M$.
\end{enumerate}
\end{problem}

\begin{problem}
Let $n \geq 3$ and let $A$ be a real $n\times n$ matrix whose entries in every row form an arithmetic progression. Prove that there exist real numbers $\alpha,\beta$ such that $A^3+\alpha A^2+\beta A = 0$.
\end{problem}

\begin{problem}
Let $A=\begin{pmatrix}1&1\\-1&-1\end{pmatrix}$, $B=\begin{pmatrix}1&-1\\1&-1\end{pmatrix}$, and let $M_1,\dots,M_{10}$ be $10$ matrices satisfying:
\begin{enumerate}[label=\roman*)]
\item all entries of every $M_i$ ($i=1,\dots,10$) are integers;
\item for each $i=1,\dots,10$, $M_i$ commutes with $A$ or with $B$;
\item $M_1^2+M_2^2+\cdots+M_{10}^2 = \begin{pmatrix}2016 & 0\\ 0 & -2008\end{pmatrix}$.
\end{enumerate}
Compute $N = M_1^6+M_2^6+\cdots+M_{10}^6$.
\end{problem}

\begin{problem}
Let $A,B$ be real $n\times n$ matrices with $A^{15}=B^{16}=I_n$ and $AB=BA$. Prove that $I_n+A+B$ is invertible.
\end{problem}

\begin{problem}
Let $N$ be an $n\times n$ matrix all of whose entries equal $1/n^2$, and let $A$ be a real $n\times n$ matrix such that $A^k = N$ for some positive integer $k$. Writing $A=[a_{ij}]_n$, prove that $\displaystyle\sum_{i,j=1}^n a_{ij}^2 \geq 1$.
\end{problem}

\begin{problem}
Let $A,B$ be $n\times n$ matrices with $A$ invertible and $A^2-B=AB$. Can $B$ be nilpotent? Justify your answer.
\end{problem}

\begin{problem}
Let $A=[a_{ij}]_{i,j=1}^{2021}$ be a nonsingular real matrix of order $2021$. Consider the column matrices $A_k=[a_{1k},a_{2k},\dots,a_{2021,k}]^t$, $k=1,\dots,2021$. Let $M=A_1+2A_2+\cdots+2021A_{2021}$. Find all matrices $X$ satisfying:
\[
\begin{aligned}
X(M-A_1) &= A_1-A_{2021},\\
X(M-A_2) &= \tfrac{1}{2}(A_2-A_{2020}),\\
&\ \ \vdots\\
X(M-A_k) &= \tfrac{1}{k}(A_k-A_{2022-k}),\\
&\ \ \vdots\\
X(M-A_{2021}) &= \tfrac{1}{2021}(A_{2021}-A_1).
\end{aligned}
\]
\end{problem}

\begin{problem}
Let $f(x)=\displaystyle\sum_{i=0}^{2021}(-1)^i(i+1)x^i$ and
\[ A=\begin{pmatrix}2&-5&4\\3&-7&5\\4&-9&6\end{pmatrix}. \]
Compute $f(A)$.
\end{problem}

\begin{problem}
Let $r$ be a positive integer and $\xi=\exp\!\left(\dfrac{2i\pi}{r}\right)$ a complex $r$-th root of unity. For $\alpha \in \mathbb{R}$, define $\xi^\alpha = \exp\!\left(\dfrac{2i\pi\alpha}{r}\right)$. Find the inverse of the matrix
\[
A = \begin{bmatrix}
1 & \xi^\alpha & \cdots & \xi^{(r-1)\alpha}\\
1 & \xi^{\alpha+1} & \cdots & \xi^{(r-1)(\alpha+1)}\\
\vdots & \vdots & & \vdots\\
1 & \xi^{\alpha+k} & \cdots & \xi^{(r-1)(\alpha+k)}\\
\vdots & \vdots & & \vdots\\
1 & \xi^{\alpha+r-1} & \cdots & \xi^{(r-1)(\alpha+r-1)}
\end{bmatrix}.
\]
\end{problem}

\begin{problem}
Let $A=(a_{ij})$ be a $2021\times 2021$ matrix with entries $a_{ij}=i^2+j^2+2020ij+2021$ for all $i,j=1,\dots,2021$. Compute $\det A$.
\end{problem}

\begin{problem}
Let $a_1,a_2,a_3,a_4$ and $b_1,b_2,b_3,b_4$ be, respectively, $4$ consecutive even numbers and $4$ consecutive odd numbers. Compute $\det A$, where
\[
A = \begin{pmatrix}
(a_1+b_1)^3 & (a_1+b_2)^3 & (a_1+b_3)^3 & (a_1+b_4)^3\\
(a_2+b_1)^3 & (a_2+b_2)^3 & (a_2+b_3)^3 & (a_2+b_4)^3\\
(a_3+b_1)^3 & (a_3+b_2)^3 & (a_3+b_3)^3 & (a_3+b_4)^3\\
(a_4+b_1)^3 & (a_4+b_2)^3 & (a_4+b_3)^3 & (a_4+b_4)^3
\end{pmatrix}.
\]
\end{problem}

\begin{problem}
Let $A=[a_{ij}]_{6\times 6}$ with $a_{ij}=0$ if $i=j$, and $a_{ij}\in\{1,2021\}$ if $i\neq j$. Prove that $\det A$ is always nonzero.
\end{problem}

\begin{problem}
Let $A=(a_{ij})$, $B=(b_{ij})$ be real $n\times n$ matrices satisfying $AB^2+A = 2AB+I$.
\begin{enumerate}[label=\alph*)]
\item Prove that $A$ and $B$ commute.
\item If $a_{ij},b_{ij}$ are integers for all $i,j\in\{1,\dots,n\}$, compute $\det A$.
\end{enumerate}
\end{problem}

\begin{problem}
Let $X$ be a real $n\times n$ matrix satisfying $X+X^T=I_n$ ($I_n$ the identity matrix, $X^T$ the transpose of $X$). Prove that $\det X \geq \dfrac{1}{2^n}$.
\end{problem}

\begin{problem}
Compute
\[
\lim_{n\to+\infty}
\begin{vmatrix}
x & 0 & 0 & \cdots & 0 & 1/n!\\
-1 & x & 0 & \cdots & 0 & 1/n!\\
0 & -2 & x & \cdots & 0 & 1/n!\\
\vdots & \vdots & \vdots & \ddots & \vdots & \vdots\\
0 & 0 & 0 & \cdots & x & 1/n!\\
0 & 0 & 0 & \cdots & -n & 1/n!
\end{vmatrix}.
\]
\end{problem}

\begin{problem}
Let $A,B\in\mathcal{M}_n(\mathbb{R})$ satisfy $AB=BA$, with $B$ nilpotent. Prove that $\det(A+B)=\det(A)$.
\end{problem}

\begin{problem}
Let $A=(a_{ij})$ be a $2022\times2022$ matrix with $a_{ij}=-a_{ji}$ for all $i,j=\overline{1,2022}$. Let $B=(b_{ij})$ be a matrix of the same order with $b_{ij}=a_{ij}+23$ for all $i,j=\overline{1,2022}$. Prove that $\det A = \det B$.
\end{problem}

\begin{problem}
Let $A$ be an $n\times n$ matrix with $\mathrm{rank}\,A=1$. Prove that
\[ \det(A+I_n) = \mathrm{tr}(A)+1. \]
\end{problem}

\begin{problem}
Let real numbers $a_{ij}+a_{ji}=0$ for $i,j=1,2,\dots,2014$. Solve the linear system
\[
\begin{cases}
(a_{11}+2014)x_1 + a_{12}x_2 + \cdots + a_{1,2014}x_{2014} = 0\\
a_{21}x_1 + (a_{22}+2014)x_2 + \cdots + a_{2,2014}x_{2014} = 0\\
\qquad\vdots\\
a_{2014,1}x_1 + a_{2014,2}x_2 + \cdots + (a_{2014,2014}+2014)x_{2014} = 0
\end{cases}
\]
\end{problem}

\begin{problem}
Let $P(x)=(x-2)(x-3)\cdots(x-2014)$. Suppose $P(x)=a_1+a_2x+\cdots+a_{2014}x^{2013}$. Solve the system
\[
\begin{cases}
a_1x_1 + a_2x_2 + \cdots + a_{2013}x_{2013} + a_{2014}x_{2014} = 0\\
a_{2014}x_1 + a_1x_2 + \cdots + a_{2012}x_{2013} + a_{2013}x_{2014} = 0\\
\qquad\vdots\\
a_2x_1 + a_3x_2 + \cdots + a_{2014}x_{2013} + a_1x_{2014} = 0
\end{cases}
\]
\end{problem}

\begin{problem}
Let $K_2(\mathbb{Z}) = \{A \mid A\in M_2(\mathbb{Z}) \text{ and } \det A = 1\}$.
\begin{enumerate}[label=\alph*.]
\item Do there exist $A,B,C\in K_2(\mathbb{Z})$ such that $A^2+B^2=C^2$?
\item Do there exist $A,B,C\in K_2(\mathbb{Z})$ such that $A^4+B^4=C^4$?
\end{enumerate}
\end{problem}

\begin{problem}
Define the trace of a square matrix $C=(c_{ij})_{n\times n}$, denoted $\mathrm{tr}(C)$, as $\displaystyle\sum_{i=1}^n c_{ii}$. Let $A,B$ be real $n\times n$ matrices satisfying $AB-BA=A$. Compute the trace of $A^{2014}$.
\end{problem}

\begin{problem}
Define the trace of a matrix $A$, denoted $\mathrm{Tr}(A)$, as the sum of all entries on the main diagonal of $A$. Let $A=(a_{ij})_{n\times n}$ be an $n\times n$ matrix satisfying $A^{2014}=0$. Compute $\mathrm{Tr}(A)$.
\end{problem}

\begin{problem}
Let $A$ be an $n\times n$ matrix with $2015$ on the main diagonal and $2014$ elsewhere:
\[
A = \begin{pmatrix}
2015 & 2014 & \cdots & 2014\\
2014 & 2015 & \cdots & 2014\\
\vdots & & \ddots & \vdots\\
2014 & 2014 & \cdots & 2015
\end{pmatrix}.
\]
Compute $A^k$, where $k$ is a given positive integer.
\end{problem}

\begin{problem}
Let $A,B,C,D$ be $n\times n$ matrices satisfying $BA^T$ and $DC^T$ are symmetric, and $DA^T-CB^T=I$. Prove that $D^TA-B^TC=I$.
\end{problem}

\begin{problem}
For each matrix $A\in M_2(\mathbb{R})$, determine whether there exist $B,C\in M_2(\mathbb{R})$ such that $A=B^2+C^2$.
\end{problem}

\begin{problem}
Let $A,B$ be two $2\times 2$ matrices satisfying $(AB)^2=0$. Prove that $(BA)^2=0$. Does the statement remain true for $A,B$ square matrices of order greater than $2$?
\end{problem}

\begin{problem}
\begin{enumerate}[label=\arabic*.]
\item Let $A=\begin{pmatrix}1&0&0\\1&1&0\\1&1&1\end{pmatrix}$ and let $n>0$ be an integer. Find $(A^n)^{-1}$.
\item Let $A$ and $B$ be distinct real $n\times n$ matrices. Suppose $A^3=B^3$ and $A^2B=B^2A$. Prove that $A^2+B^2$ is not invertible.
\end{enumerate}
\end{problem}

\begin{problem}
Compute $\det A$, where $A$ is the $2000\times 2000$ matrix with $(i,j)$-entry equal to $|i-j|+1$.
\end{problem}

\begin{problem}
Let $a_0$ be a real number, and suppose $\{a_0,a_1,a_2,\dots,a_{2013}\}$ is an arithmetic progression with common difference $d=4$. Find the condition on $a_0$ such that the symmetric (Toeplitz-type) matrix $A=(a_{|i-j|})_{0\le i,j\le 2013}$ is invertible.
\end{problem}

\begin{problem}
Let $A$ and $B$ be two $2013\times 2013$ matrices satisfying $AB^2A+BA^2B=I$, where $I$ is the $2013\times 2013$ identity matrix. Find the sum of the entries on the main diagonal of $AB^2A$.
\end{problem}

\begin{problem}
Let $A$ and $B$ be two $2013\times 2013$ matrices satisfying $\mathrm{rank}(AB)=\mathrm{rank}(A)\cdot\mathrm{rank}(B)$. Determine the matrix $A$ if $\mathrm{rank}(B)\geq 2$.
\end{problem}

\begin{problem}
Let $A$ be a $3\times 3$ matrix whose entries are $0$ or $1$. Find the maximum value of $\det(A)$.
\end{problem}

\begin{problem}
Does there exist a $3\times 3$ square matrix $A$ such that $\mathrm{Tr}(A)=0$ and $A^T+A^2=I$, where $\mathrm{Tr}(A)$ is the sum of the entries on the main diagonal of $A$?
\end{problem}

\begin{problem}
Let $A,B$ be $n\times n$ matrices such that $A^{2013}=0$, $AB=BA$, $B\neq 0$. Prove that $\mathrm{rank}(AB) \leq \mathrm{rank}(B)-1$.
\end{problem}

\begin{problem}
\begin{enumerate}[label=\arabic*.]
\item Find all matrices commuting with $A=\begin{pmatrix}0&1&2\\0&0&3\\0&0&0\end{pmatrix}$.
\item Solve the equation $X^n=A$ for $n\in\mathbb{N}^*$.
\end{enumerate}
\end{problem}

\begin{problem}
Let $a,b,c$ be real numbers satisfying $a^2+b^2+c^2=4$. Find the maximum and minimum of the determinant of
\[
A = \begin{pmatrix}
a+b & b+c & c+a\\
c+a & a+b & b+c\\
b+c & c+a & a+b
\end{pmatrix}.
\]
\end{problem}

\begin{problem}
Compute the sum of all determinants of $n\times n$ matrices ($n\geq 2$) such that every row and every column contains exactly one nonzero entry, and the nonzero entries are pairwise distinct values taken from $\{1,2,\dots,n\}$.
\end{problem}

\begin{problem}
Suppose $A$ is a $2013\times 2013$ matrix satisfying: $\mathrm{trace}(A^2)=8052$, and every $2013\times 2013$ matrix $B$ can be written as $B=B_1+B_2$, where $AB_1=B_1A$ and $AB_2=-B_2A$. Prove that there exists a natural number $m\leq 2013$ such that $\det(A-I)=(-3)^m$.
\end{problem}

\begin{problem}
Do there exist two $2\times 2$ matrices $A,B$ such that $C=AB-BA$ commutes with both $A$ and $B$, and $C$ is nonzero?
\end{problem}

\begin{problem}
Let $A,B$ be $n\times n$ matrices satisfying $\mathrm{Im}(A)\cap\mathrm{Im}(B)=\{0\}$, and let $\{u_1,\dots,u_k\}$, $\{v_1,\dots,v_k\}$ be arbitrary subsets of $\mathbb{R}^n$. Prove that if $k > r(A)+r(B)$ (where $r(A)$ denotes the rank of $A$) then there always exist real numbers $\lambda_1,\dots,\lambda_k$, not all zero, such that
\[ \lambda_1Au_1+\lambda_2Au_2+\cdots+\lambda_kAu_k = \lambda_1Bv_1+\lambda_2Bv_2+\cdots+\lambda_kBv_k. \]
\end{problem}

\begin{problem}
Let $V$ be the set of $n\times n$ matrices whose entries are pairwise distinct values taken from $\{1,2,\dots,n^2\}$. Find the maximum and minimum of $r(A)$ for $A\in V$ ($r(A)$ denotes the rank of $A$).
\end{problem}

\begin{problem}
Find all $n\times n$ matrices $A$ such that for every $n\times n$ matrix $B$, $\det(A+2013B)=\det A + 2013\cdot\det B$.
\end{problem}

\begin{problem}
\begin{enumerate}[label=\arabic*.]
\item Let $A,B\in \mathrm{Mat}(n,\mathbb{R})$ be such that there exist $(\alpha,\beta)\in(\mathbb{R}-\{0\})^2$ satisfying $AB+\alpha A+\beta B=0$. Prove that $AB=BA$.
\item Prove that for all $A,B,C\in \mathrm{Mat}(2,\mathbb{R})$,
\[ (AB-BA)^2C - C(AB-BA)^2 = O. \]
\end{enumerate}
\end{problem}

\begin{problem}
Suppose the determinant of the matrix $A=[a_{ij}]_{n\times n}$ equals $\alpha$, and the sum of the algebraic cofactors of the entries of $A$ equals $\beta$ ($\alpha,\beta\in\mathbb{R}$). Compute the determinants of the following matrices:
\begin{enumerate}[label=\arabic*.]
\item $B = [a_{ij}+2013]_{n\times n}$ (every entry of $A$ increased by $2013$).
\item $C$, the matrix whose first row consists of all $1$'s and whose $(k+1)$-th row (for $k=1,\dots,n-1$) is $(a_{k+1,1}-a_{11},\, a_{k+1,2}-a_{12},\,\dots,\, a_{k+1,n}-a_{1n})$.
\end{enumerate}
\end{problem}

\begin{problem}
Let $A$ and $B$ be $n\times n$ matrices satisfying $\mathrm{rank}(AB)=\mathrm{rank}(B)$. Prove that $ABX=ABY \iff BX=BY$ for all $X,Y$.
\end{problem}

\begin{problem}
Let $A$ and $B$ be $n\times n$ orthogonal matrices satisfying $\det(AB)<0$. Prove that $\det A + \det B = \det(A+B)$.
\end{problem}

\begin{problem}
Let $A$ be an $n\times n$ matrix. Prove that if $\mathrm{trace}(A^TA)+n = 2\cdot\mathrm{trace}(A)$ then $A$ is invertible.
\end{problem}

\begin{problem}
Let $A$ and $B$ be two $2013\times 2013$ matrices satisfying $AB+2012A+2013B=0$. Prove that $\mathrm{rank}(A)+\mathrm{rank}(B) \neq 2013$.
\end{problem}

\begin{problem}
Prove that if an $n\times n$ square matrix $A$ has entries equal to $0$ on the main diagonal, and all remaining entries equal to $1$ or $2014$, then $\mathrm{rank}(A) \geq n-1$.
\end{problem}

\begin{problem}
Let $A,B$ be real square matrices satisfying $A^{2013}=0$, $AB=2012A+2011B$. Prove that $B^{2013}=0$ and $\det(A-2011I)\neq 0$.
\end{problem}

\begin{problem}
Let $f:V\to W$ be a linear map between finite-dimensional vector spaces over a field $K$. Prove that:
\begin{enumerate}[label=\arabic*.]
\item If $A$ is a $k$-dimensional subspace of $V$ such that $A\cap\mathrm{Ker}\,f$ is an $r$-dimensional subspace, then $\dim f(A) = k-r$.
\item If $B$ is a subspace of $W$ such that $B\cap \mathrm{Im}\,f$ is an $s$-dimensional subspace, then $\dim f^{-1}(B) = \dim V + s - \mathrm{rank}(f)$.
\end{enumerate}
\end{problem}

\begin{problem}
Let $V=\mathbb{F}[x]$ and let $f$ be the endomorphism of $V$ defined by $f(P)=xP$. Determine the eigenvalues and eigenvectors of the endomorphism $F:\mathrm{End}(V)\to\mathrm{End}(V)$ defined by $F(g)=f\circ g - g\circ f$.
\end{problem}

\begin{problem}
Let $A$ be a real $n\times n$ matrix and let $\varphi_A,\psi_A$ be the linear endomorphisms of the real vector space $M(n,\mathbb{R})$ of $n\times n$ real matrices, defined by $\varphi_A(X)=AX-XA$, $\psi_A(X)=AX$. Prove that $\det(\varphi_A)=0$ and $\det(\psi_A)=(\det A)^n$.
\end{problem}

\begin{problem}
A polynomial in $n$ variables $x_1,\dots,x_n$ is called \emph{homogeneous of degree $d$} if the sum of exponents of the variables in every monomial equals $d$ (for example, $x_1^2x_2^3+x_3^5$ is homogeneous of degree $5$ in three variables $x_1,x_2,x_3$). Find the dimension of the real vector space consisting of the zero polynomial together with all homogeneous polynomials of degree $d$ in $n$ variables with real coefficients.
\end{problem}

\begin{problem}
Let $V$ be an $n$-dimensional vector space over $\mathbb{R}$ and let $f$ be a linear endomorphism of $V$. For each nonnegative integer $k$, let $f^k$ denote the $k$-fold composition of $f$ with itself ($f^0=\mathrm{Id}_V$, $f^1=f$, $f^2=f\circ f$, \dots). Prove that:
\begin{enumerate}[label=\arabic*.]
\item There exists a nonnegative integer $m\leq n$ such that
\[
\mathrm{Im}\,f^0 \supsetneq \mathrm{Im}\,f^1 \supsetneq \cdots \supsetneq \mathrm{Im}\,f^m = \mathrm{Im}\,f^{m+1} = \cdots,
\]
\[
\ker f^0 \subsetneq \ker f^1 \subsetneq \cdots \subsetneq \ker f^m = \ker f^{m+1} = \cdots.
\]
\item $\{\dim\ker f^{k+1}-\dim\ker f^k\}_{k=0,1,\dots}$ is a decreasing sequence.
\item $V = \ker f^m \oplus \mathrm{Im}\,f^m$.
\item Every $n\times n$ matrix with real coefficients is equivalent to a matrix of the form $\begin{pmatrix}N & 0\\0 & C\end{pmatrix}$, where $N$ is a nilpotent square matrix and $C$ is an invertible square matrix.
\end{enumerate}
\end{problem}

\begin{problem}
Let $X$ be a set. Equip $F(X,\mathbb{R})$ with the $\mathbb{R}$-vector space structure induced from that of $\mathbb{R}$. Let $n\geq 1$ and $f_1,\dots,f_n\in F(X,\mathbb{R})$. Prove that $f_1,\dots,f_n$ are linearly independent if and only if there exist $n$ elements $x_1,\dots,x_n\in X$ such that the matrix $(f_i(x_j))_{1\leq i,j\leq n}$ is invertible.
\end{problem}

\begin{problem}
\begin{enumerate}[label=\arabic*.]
\item Let $f$ be a linear endomorphism of $\mathbb{C}^n$. Prove that the following two conditions are equivalent: (i) $f$ is diagonalizable; (ii) $f^2$ is diagonalizable and $\ker f = \ker f^2$.
\item Find necessary and sufficient conditions on the tuple of complex numbers $(a_1,\dots,a_n)$ so that the matrix
\[
\begin{pmatrix}
0 & 0 & \cdots & a_n\\
0 & 0 & \cdots & 0\\
\vdots & \vdots & \cdots & \vdots\\
0 & a_2 & \cdots & 0\\
a_1 & 0 & \cdots & 0
\end{pmatrix}
\]
is diagonalizable.
\end{enumerate}
\end{problem}

\begin{problem}
Let $V$ be a vector space over a field $K$. Let $F_1,F_2,\dots,F_n$ ($n\in\mathbb{N}, n\geq 1$) be finitely generated subspaces of $V$. Prove that
\[ \dim(F_1+F_2+\cdots+F_n) \leq \dim F_1+\dim F_2+\cdots+\dim F_n. \]
Moreover, equality holds if and only if $F_1+F_2+\cdots+F_n$ is a direct sum in $V$.
\end{problem}

\begin{problem}
Let $V_1,V_2,V_3$ be three subspaces of a $2014$-dimensional vector space $V$ over $\mathbb{R}$. Prove that if $\dim V_1+\dim V_2+\dim V_3 > 4028$ then $V_1\cap V_2\cap V_3 \neq 0$.
\end{problem}

\begin{problem}
Let $E$ be a finite-dimensional $\mathbb{R}$-vector space, and call an endomorphism $u$ of $E$ nilpotent if there is a positive integer $k$ with $u^k=0$. Define the nilpotency degree of a nilpotent endomorphism $u$ as the smallest positive integer $p$ such that $u^p=0$. Let $u$ be a nilpotent endomorphism of degree $p$. Prove that there exists a vector $x$ of $E$ such that the system $\{x,u(x),\dots,u^{p-1}(x)\}$ is linearly independent.
\end{problem}

\begin{problem}
Let $E$ be a vector space over $\mathbb{R}$. A linear endomorphism $f$ of $E$ is called nilpotent if there exists $k\in\mathbb{N}$ such that $f^k=0$ (the smallest such $k$ is called the nilpotency degree of $f$). Prove that if $f$ is a nilpotent endomorphism of degree $p$ of $E$, then the system $\{\mathrm{id}_E, f,\dots, f^{p-1}\}$ is linearly independent in the vector space $\mathrm{End}(E)$.
\end{problem}

\begin{problem}
Let $V$ be the space of continuous real functions on $[a,b]$, and suppose $f_1,f_2,\dots,f_n\in V$ satisfy $\det\left(\int_a^b f_i(x)f_j(x)\,dx\right)=0$. Prove that the system of vectors $\{f_1,f_2,\dots,f_n\}$ is linearly dependent in $V$.
\end{problem}

\begin{problem}
Let $A,B$ be $n\times n$ matrices, where $B$ is nilpotent and $AB=BA$. Prove that $\lambda$ is an eigenvalue of $A$ if and only if $\lambda$ is an eigenvalue of $A+\lambda B$.
\end{problem}

\begin{problem}
Let $A=(a_{ij})$ be an $n\times n$ matrix, where $a_{ij}=b^{i-j}$ if $i\neq j$, and $a_{ij}=0$ if $i=j$, with $b\neq 0$. Find the eigenvalues of $A$ and compute $\det(A)$.
\end{problem}

\begin{problem}
Let $A,B$ be $n\times n$ matrices satisfying $A-A^2+B-B^2=AB$. Prove that $A+B$ has eigenvalue $1$ if and only if $A$ or $B$ has eigenvalue $1$.
\end{problem}

\begin{problem}
Let $A,B$ be real square matrices satisfying $A^TA=I$ and $B^TB=I$, with $\det A \neq \det B$. Prove that $\det(A+B)=0$.
\end{problem}

\begin{problem}
Let $A,B,C$ be pairwise commuting real square matrices. Prove that there exist real numbers $a,b,c$, not all zero, such that $\det(aA+bB+cC)=0$.
\end{problem}

\begin{problem}
Let $A=[a_{ij}]$ be an $n\times n$ matrix defined by $a_{ij}=0$ if $i^j+j^i=3k$; $a_{ij}=1$ if $i^j+j^i=3k+1$; $a_{ij}=2$ if $i^j+j^i=3k+2$. Find the largest natural number $n$ such that $\det(A)\neq 0$.
\end{problem}

\begin{problem}
Compute the determinant
\[
\begin{vmatrix}
0 & 1 & 2^{2014} & 3^{2014} & \cdots & 2014^{2014}\\
1 & 0 & 1 & 2^{2014} & \cdots & 2013^{2014}\\
2^{2014} & 1 & 0 & 1 & \cdots & 2012^{2014}\\
\vdots & & & & \ddots & \vdots\\
2014^{2014} & 2013^{2014} & 2012^{2004} & 2011^{2014} & \cdots & 0
\end{vmatrix}.
\]
\end{problem}

\begin{problem}
Let $A\in M_{2014\times n}(\mathbb{R})$, $n\geq 1$. Prove that $\det(AA^T+4I_{2014}) \geq 2^{4028}$.
\end{problem}

\begin{problem}
Let $A\in M_2(\mathbb{R})$, $B\in M_3(\mathbb{R})$. Prove that
\begin{enumerate}[label=\alph*.]
\item $\det(A) = \dfrac12\left[(\mathrm{tr}(A))^2-\mathrm{tr}(A^2)\right]$.
\item $\det(B) = \dfrac16\left[(\mathrm{tr}(B))^3-3\,\mathrm{tr}(B)\,\mathrm{tr}(B^2)+2\,\mathrm{tr}(B^3)\right]$.
\end{enumerate}
\end{problem}

\begin{problem}
Compute the determinant
\[
\begin{vmatrix}
a_1 & 1 & \cdots & 1 & 1\\
-1 & a_2 & \cdots & 1 & 1\\
\vdots & \vdots & \vdots & \vdots & \vdots\\
-1 & -1 & \cdots & \cdots & 1\\
-1 & -1 & \cdots & -1 & a_{2014}
\end{vmatrix}.
\]
\end{problem}

\begin{problem}
Let $A=(a_{ij})$ be an $n\times n$ matrix ($n\geq 2$) such that every entry $a_{ij}$ has the form $p^2$, where $p$ is a prime greater than $3$. Prove that $\det(A)$ is divisible by $24^{n-1}$.
\end{problem}

\begin{problem}
Let $A$ be a square matrix of order $n>1$: $A=(a_{ij})$, $a_{ij}\in\mathbb{Z}$, where $a_{ij}$ is odd for $i\neq j$ and $a_{ii}$ is even ($1\leq i,j\leq n$). Prove that $\det(A)\neq 0$.
\end{problem}

\begin{problem}
Let $A$ be a $3\times 3$ matrix: $A=(a_{ij})$ with $a_{ij}\in K$, $1\leq i,j\leq n$, $K$ a field. Prove that $A^2=0$ if and only if $\mathrm{rank}(A)\leq 1$ and $\mathrm{trace}(A)=0$.
\end{problem}

\begin{problem}
Compute the determinant
\[
D = \begin{vmatrix}
x & 1 & 0 & 0 & \cdots & 0 & 0\\
n-1 & x & 2 & 0 & \cdots & 0 & 0\\
0 & n-2 & x & 3 & \cdots & 0 & 0\\
\vdots & & & \ddots & & & \vdots\\
0 & 0 & 0 & 0 & \cdots & x & n-1\\
0 & 0 & 0 & 0 & \cdots & 1 & x
\end{vmatrix}.
\]
\end{problem}

\begin{problem}
Let $A$ be a $2013\times 2013$ matrix. Prove that if $\det(A^{-1})=2013$ then the entries of $A$ cannot all be integers.
\end{problem}

\begin{problem}
Let $f(x)=a_1+a_2x+\cdots+a_nx^{n-1}$ be a polynomial with real coefficients, and let $\omega_1,\omega_2,\dots,\omega_n$ be the $n$-th roots of unity. Let $A$ be the circulant matrix whose $k$-th row is a cyclic shift of $(a_1,\dots,a_n)$, and let $B$ be the Vandermonde-type matrix with $(i,j)$-entry $\omega_j^{\,i-1}$. Compute $\det(A)$.
\end{problem}

\begin{problem}
Let $W$ be the set of $3\times 3$ matrices whose entries take only the values $\pm 1$. Find the number of matrices in $W$ with positive determinant.
\end{problem}

\begin{problem}
a) Prove that
\[
\det\begin{pmatrix}
1 & a_1 & a_1(a_1-1) & a_1(a_1-1)(a_1-2)\\
1 & a_2 & a_2(a_2-1) & a_2(a_2-1)(a_2-2)\\
1 & a_3 & a_3(a_3-1) & a_3(a_3-1)(a_3-2)\\
1 & a_4 & a_4(a_4-1) & a_4(a_4-1)(a_4-2)
\end{pmatrix}
= \prod_{1\leq i<j\leq 4}(a_j-a_i).
\]
b) Assuming $a_1,a_2,a_3,a_4$ are integers, prove that $\prod_{1\leq i<j\leq 4}(a_j-a_i)$ is divisible by $12$.
\end{problem}

\begin{problem}
Let $a_1,a_2,a_3$ be distinct real numbers. Prove that for every triple of real numbers $b_1,b_2,b_3$, there exists a unique polynomial $P(x)$ of degree at most $5$ satisfying $P(a_i)=P'(a_i)=b_i$, $i=1,2,3$, where $P'$ denotes the derivative of $P$.
\end{problem}

\begin{problem}
a) Let $V_4$ denote the vector space of polynomials with real coefficients of degree at most $4$. Define $e:V_4\to V_4$ by $e(f) := \displaystyle\sum_{i=0}^4 \frac{f^{(i)}}{i!}$, where $f^{(i)}$ is the $i$-th derivative of $f$ ($f^{(0)}=f$). Prove that $e$ is an invertible linear map from $V_4$ to itself. b) Let $V$ denote the vector space of polynomials with real coefficients. For each polynomial $f$, define $e(f) := \displaystyle\sum_{i=0}^\infty \frac{f^{(i)}}{i!}$. Prove that $e$ is an invertible linear map from $V$ to itself.
\end{problem}

\begin{problem}
a) Let $X = \begin{pmatrix}E_m & B\\ C & E_n\end{pmatrix}$ be a block matrix formed from the identity matrices $E_m,E_n$ of orders $m,n$ respectively, and matrices $B,C$ of sizes $m\times n$ and $n\times m$ respectively. Prove that
\[ \det(X) = \det(E_n-CB) = \det(E_m-BC). \]
b) More generally, let $X=\begin{pmatrix}A & B\\ C & D\end{pmatrix}$, where $A,D$ are square matrices and $A$ is invertible. Prove that $\det(X)=\det(A)\det(D-CA^{-1}B)$.
\end{problem}

\begin{problem}
Let $A$ and $B$ be square matrices of order $2010$ with real coefficients such that
\[ \det A = \det(A+B) = \det(A+2B) = \cdots = \det(A+2010B) = 0. \]
a) Prove that $\det(xA+yB)=0$ for all $x,y\in\mathbb{R}$. b) Give an example showing that the conclusion of part a) is no longer true if one only assumes $\det A = \det(A+B) = \cdots = \det(A+2009B)=0$.
\end{problem}

\begin{problem}
a) Let $A$ be a real square matrix of order $n\geq 2$ with the sum of its diagonal entries equal to $10$ and $\mathrm{rank}(A)=1$. Find the characteristic polynomial and the minimal polynomial of $A$ (the minimal polynomial of $A$ is the polynomial $P(t)=0$ of smallest degree, with real coefficients and leading coefficient $1$, such that $P(A)=0$). b) Let $A,B,C$ be real square matrices of order $n$, where $A$ is invertible and commutes with both $B$ and $C$. Suppose $C(A+B)=BC$. Prove that $B$ and $C$ commute with each other.
\end{problem}

\begin{problem}
Solve the system of linear equations
\[
\begin{cases}
x_1 = 2(4x_1+3x_2+2x_3+x_4)\\
x_2 = 3(x_1+4x_2+3x_3+2x_4)\\
x_3 = 4(2x_1+x_2+4x_3+3x_4)\\
x_4 = 5(3x_1+2x_2+x_3+4x_4)
\end{cases}
\]
\end{problem}

\begin{problem}
A square matrix $A$ is called \emph{nilpotent} if there exists $k>0$ such that $A^k=0$.
\begin{enumerate}[label=\alph*)]
\item Prove that every strictly upper triangular $n\times n$ matrix is nilpotent, and that the set of such matrices forms a subspace $V_0$ of $M_n(\mathbb{R})$. Compute $\dim V_0$.
\item Suppose $V$ is a subspace of $M_n(\mathbb{R})$, all of whose elements are nilpotent matrices. Prove that $\dim V \leq \dfrac{n^2-n}{2}$.
\end{enumerate}
\end{problem}

\begin{problem}
Let $A$ be a square matrix of order $n\geq 2$ whose entries are odd perfect squares. Prove that $\det(A)$ is divisible by $8^{n-1}$.
\end{problem}

\begin{problem}
Let $A,B\in M_{100}(\mathbb{R})$ be two matrices satisfying $A^{101}=0$ and $AB=2A+3B$. Prove that $(A+B)^{100}=0$.
\end{problem}

\begin{problem}
Let $A$ be a square matrix of order $5$ whose entries are $1$ or $-1$. Prove that $|\det(A)|\leq 64$.
\end{problem}

\begin{problem}
Let $A \in \mathcal{M}_n(\mathbb{C})$ be a matrix of rank $1$ such that $\operatorname{Tr} A \in \mathbb{R}$ and there exist $a,b\in\mathbb{C}$ with $|a|=1$ such that
\[
A^2 + aAA^* + a A^*A + b(A^*)^2 = O_n,
\]
where $A^* = \overline{A}^{\,T}$. Prove that there exists $u \in \mathbb{C}^n$ such that $A = \pm u\cdot \bar u^T$.
\end{problem}

\begin{problem}
Let $A, B \in \mathcal{M}_n(\mathbb{R})$ such that $A^2 = -I_n$, $\det B \ne 0$, and $AB = -BA$. Prove that $n$ is even and the sign of $\det B$ is $(-1)^{n/2}$.
\end{problem}

\begin{problem}
Let $X, Y \in \mathcal{M}_n(\mathbb{C})$ such that $Y^2 = YX - XY$ and the rank of $X+Y$ is $1$. Prove that $Y^3 = YXY = O_n$.
\end{problem}

\begin{problem}
Consider $n \ge 2$ and $A \in \mathcal{M}_n(\mathbb{C})$. Prove that there exists $k \in \mathbb{N}$, $1 \le k \le n$, and $B \in \mathcal{M}_n(\mathbb{C})$ such that
\[
B^2 = B, \qquad \ker B = \ker A^k, \qquad \operatorname{Im} B = \operatorname{Im} A^k.
\]
\end{problem}

\begin{problem}
Let $A, B \in \mathcal{M}_n(\mathbb{C})$. Prove that the following statements are equivalent:
\begin{enumerate}[label=\arabic*)]
\item For all matrices $C \in \mathcal{M}_n(\mathbb{C})$, there exist $X,Y \in \mathcal{M}_n(\mathbb{C})$ such that $AX+YB=C$;
\item For all matrices $C \in \mathcal{M}_n(\mathbb{C})$, there exist $U,V \in \mathcal{M}_n(\mathbb{C})$ such that $A^2U+VB^2=C$.
\end{enumerate}
\end{problem}

\begin{problem}
Consider $A, B \in \mathcal{M}_n(\mathbb{C})$ such that $AB+BA = A+B$. Prove that the following assertions are equivalent:
\begin{enumerate}[label=\alph*)]
\item $\operatorname{rank}AB + n = \operatorname{rank}A + \operatorname{rank}B$;
\item $\ker B \subset \operatorname{Im}A$.
\end{enumerate}
\end{problem}

\begin{problem}
Let $A,B,C \in \mathcal{M}_n(\mathbb{C})$ with $A \ne O$ such that $A + BAC = BA + AC$. Prove that $1$ is an eigenvalue either for matrix $B$ or for matrix $C$, with algebraic multiplicity at least $\dfrac{\operatorname{rank}A}{2}$.
\end{problem}

\begin{problem}
We consider a natural number $n$, $n \ge 2$, and the matrix
\[
A =
\begin{pmatrix}
1 & 2 & 3 & \dots & n \\
n & 1 & 2 & \dots & n-1 \\
n-1 & n & 1 & \dots & n-2 \\
\vdots & \vdots & \vdots & \ddots & \vdots \\
2 & 3 & 4 & \dots & 1
\end{pmatrix}.
\]
Show that
\[
\varepsilon^n \det\!\left(I_n - A^{2n}\right) + \varepsilon^{n-1}\det\!\left(\varepsilon I_n - A^{2n}\right) + \varepsilon^{n-2}\det\!\left(\varepsilon^2 I_n - A^{2n}\right) + \cdots + \det\!\left(\varepsilon^n I_n - A^{2n}\right)
\]
\[
= n(-1)^{n-1}\left[\frac{n^n(n+1)}{2}\right]^{2n^2-4n}\left(1+(n+1)^{2n}\left(2n+(-1)^n\binom{2n}{n}\right)\right),
\]
where $\varepsilon \in \mathbb{C}\setminus\mathbb{R}$, $\varepsilon^{n+1}=1$.
\end{problem}

\begin{problem}
Prove that if $A \in \mathcal{M}_n(\mathbb{C})$ is invertible and anti-Hermitian, then the function
\[
f : \mathcal{M}_n(\mathbb{C}) \to \mathcal{M}_n(\mathbb{C}), \qquad f(X) = AX - XA^2, \qquad X \in \mathcal{M}_n(\mathbb{C}),
\]
is bijective.
\end{problem}

\begin{problem}
We consider a prime number $n \ge 3$, two matrices $A, B \in \mathcal{M}_{n-1}(\mathbb{Q})$, $AB = BA$, and
\[
\varepsilon = \cos\frac{2(n-1)\pi}{n} + i\sin\frac{2(n-1)\pi}{n},
\]
such that
\begin{enumerate}[label=\alph*).]
\item $\det(A - \varepsilon B) = 0$;
\item $\displaystyle\sum_{k=0}^{n-1}\varepsilon^{n-ki}\det(1-\varepsilon^{-k})\det(\varepsilon^k B - A) \ge 0$ for all $i \in \{0,1,\dots,n-2\}$;
\item $\det B > 0$.
\end{enumerate}
Show that whatever $m \ge 1$ is a natural number and whatever the rational numbers $b_0, b_1, \dots, b_m$, if $\displaystyle\sum_{k=0}^m b_k A^k B^{m-k} \ne O_n$, then
\[
\det\!\left(\sum_{k=0}^m b_k A^k B^{m-k}\right) \ne 0.
\]
\end{problem}

\begin{problem}
Let $A, B \in \mathcal{M}_n(\mathbb{R})$ be two real, symmetric matrices with nonnegative eigenvalues. Prove that
\[
A^3 + B^3 = (A+B)^3 \quad \text{if and only if} \quad AB = O_n.
\]
\end{problem}

\begin{problem}
Let $A \in \mathcal{M}_n(\mathbb{R})$ be a matrix with real eigenvalues such that
\[
\operatorname{Tr}(A^2) = \operatorname{Tr}(A^4) = \operatorname{Tr}(A^6).
\]
Prove that $\operatorname{Tr}(A^{2020}) = \operatorname{Tr}(A^2)$ and $\operatorname{Tr}(A^{2021}) = \operatorname{Tr}(A^3)$.
\end{problem}

\begin{problem}
Consider $C \in \mathcal{M}_n(\mathbb{C})$, $C = A + iB$, $A, B \in \mathcal{M}_n(\mathbb{R})$, with $\operatorname{rank} C = 1$. Prove that:
\begin{enumerate}[label=\alph*)]
\item $\operatorname{Tr}(A^2+B^2) \ge 0$;
\item if $\operatorname{Tr}(A^2+B^2) = 0$, then $A^2+B^2 = O_n$.
\end{enumerate}
\end{problem}

\begin{problem}
Let $n, k \in \mathbb{N}^*$, $a \in \mathbb{C}^*$ and $A \in \mathcal{M}_n(\mathbb{C})$ such that $A^{k+1} = a \cdot A$. Prove that $\operatorname{Tr}(A^k) = a \cdot \operatorname{rank} A$.
\end{problem}

\begin{problem}
For a matrix $M \in \mathcal{M}_n(\mathbb{C})$, we denote $\ker M = \{X \in \mathcal{M}_{n,1}(\mathbb{C}) : MX = O_{n,1}\}$.

Prove that, for any $A, B, C \in \mathcal{M}_n(\mathbb{C})$, the following assertions are equivalent:
\begin{enumerate}[label=(\roman*)]
\item $AB = A$, $BC = B$, $CA = C$;
\item $A^2 = A$, $B^2 = B$, $C^2 = C$, $\ker A = \ker B = \ker C$.
\end{enumerate}
\end{problem}

\begin{problem}
Let $A, B \in \mathcal{M}_n(\mathbb{C})$ be two matrices of the same rank and let $k \in \mathbb{N}^*$. Then
\[
A^{k+1}B^k = A \quad \text{if and only if} \quad B^{k+1}A^k = B.
\]
\end{problem}

\begin{problem}
Let \(A_n\) be the \(n\times n\) determinant
\[
A_n=
\begin{vmatrix}
a&b&b&\cdots&b&b\\
c&a&b&\cdots&b&b\\
c&c&a&\cdots&b&b\\
\vdots&\vdots&\vdots&\ddots&\vdots&\vdots\\
c&c&c&\cdots&c&a
\end{vmatrix},
\qquad a,b,c\in\mathbb R.
\]
Find a recurrence formula for \(A_n\).
\end{problem}

\begin{problem}
Let \(m,n\in\mathbb N^+\). Consider matrices
\[
A,B\in\mathcal M_n(\mathbb C)
\]
and complex numbers \(a_0,a_1,\ldots,a_m\), with \(a_0\ne0\) and \(a_m\ne0\), such that
\[
A\left(a_0I_n+a_1B+a_2B^2+\cdots+a_mB^m\right)=B.
\]
Prove that
\[
AB=BA.
\]
\end{problem}

\begin{problem}
Let \(A\) and \(B\) be two \(n\times n\) matrices that are orthogonal projections, i.e.
\[
A^2=A=A^*,\qquad B^2=B=B^*.
\]
Let \(\sqrt{A+B}\) denote the positive semidefinite square root of \(A+B\).
Prove that
\[
\tr(A+B)-(2-\sqrt2)\rank(AB)
\le \tr\sqrt{A+B}
\le
(\sqrt2-1)\tr(A+B)+(2-\sqrt2)\rank(A+B),
\]
and show that equality holds simultaneously if and only if
\[
AB=BA.
\]
\end{problem}

\begin{problem}
Let
\[
M=\{A_1,A_2,\ldots,A_n\}
\]
be a set of \(n\ge2\) matrices in
\[
\mathcal M_{2026}(\mathbb C)
\]
such that
\[
AB\in M\qquad\text{for all }A,B\in M.
\]
Prove that there exist
\[
\varepsilon_1,\varepsilon_2,\ldots,\varepsilon_n\in\{-1,0,1\},
\]
not all zero, such that
\[
\rank\left(\varepsilon_1A_1+\varepsilon_2A_2+\cdots+\varepsilon_nA_n\right)
\le1013.
\]
\end{problem}

\begin{problem}
Let \(A\in\mathcal M_2(\mathbb R)\) be a matrix with real entries such that
\[
\det(A^{2014}-I_2)=\det(A^{2014}+I_2)
\]
and
\[
\det(A^{2016}-I_2)=\det(A^{2016}+I_2).
\]
Prove that
\[
\det(A^n-I_2)=\det(A^n+I_2)
\]
for every \(n\in\mathbb N\).

Here \(\mathbb N\) denotes the set of positive integers and \(I_2\) is the \(2\times2\) identity matrix.
\end{problem}

\begin{problem}
Consider \(A,B\in\mathcal M_n(\mathbb C)\) such that
\[
AB=BA,\qquad \det B\ne0.
\]
\begin{enumerate}[label=\textup{(\alph*)}, leftmargin=2em]
    \item If
    \[
    |\det(A+zB)|=1
    \qquad\text{for all }z\in\mathbb C\text{ with }|z|=1,
    \]
    prove that
    \[
    A^n=0_n.
    \]
    \item Is the conclusion still true if the commutativity condition is dropped?
\end{enumerate}
\end{problem}

\begin{problem}
Let $S$ denote the vector space of real $n\times n$ skew-symmetric matrices. For a nonsingular matrix $A$, compute the determinant of the linear map
\[
T_A : S \to S, \qquad T_A(X) = AXA^t.
\]
\end{problem}

\begin{problem}
Let $M_{7\times 7}$ denote the vector space of real $7\times 7$ matrices. Let $A$ be a diagonal matrix in $M_{7\times 7}$ that has $+1$ in four diagonal positions and $-1$ in three diagonal positions. Define the linear transformation $T$ on $M_{7\times 7}$ by $T(X) = AX - XA$. What is the dimension of the range of $T$?
\end{problem}

\begin{problem}
Let $F$ be a field. For $m$ and $n$ positive integers, let $M_{m\times n}$ be the vector space of $m\times n$ matrices over $F$. Fix $m$ and $n$, and fix matrices $A$ and $B$ in $M_{m\times n}$. Define the linear transformation $T$ from $M_{n\times m}$ to $M_{m\times n}$ by
\[
T(X) = AXB.
\]
Prove that if $m \neq n$, then $T$ is not invertible.
\end{problem}

\begin{problem}
Determine all solutions to the following infinite system of linear equations in the infinitely many unknowns $x_1, x_2, \ldots$:
\[
\begin{aligned}
x_1 + x_3 + x_5 &= 0 \\
x_2 + x_4 + x_6 &= 0 \\
x_3 + x_5 + x_7 &= 0 \\
&\;\vdots
\end{aligned}
\]
How many free parameters are required?
\end{problem}

\begin{problem}
If a finite homogeneous system of linear equations with rational coefficients has a nontrivial complex solution, need it have a nontrivial rational solution? Give a proof or a counterexample.
\end{problem}

\begin{problem}
Let $A$ be a real $m \times n$ matrix with rational entries and let $b$ be an $m$-tuple of rational numbers. Assume that the system of equations $Ax = b$ has a solution $x$ in complex $n$-space $\mathbb{C}^n$. Show that the equation has a solution vector with rational components, or give a counterexample.
\end{problem}

\begin{problem}
Prove that a real $n\times n$ matrix $A$ satisfying $(A+A^T)^{k+2}+2(A+A^T)^k = 3(A+A^T)^{k+1}$ for some positive integer $k$ has nonnegative sum of entries.
\end{problem}

\begin{problem}
Let $A, B$ be nilpotent $n\times n$ complex matrices with $A^2B - 2ABA + BA^2 = A$. Show that there exists a finite set $S \subset M_n(\mathbb{C})$ containing $B$ such that $AX-XA+X+A \in S$ for all $X \in S$ if and only if $A = O_n$.
\end{problem}

\begin{problem}
Let $A \in \mathcal{M}_n(\mathbb{C})$ and $a \in \mathbb{C}$, $a \neq 0$, such that $A - A^* = 2aI_n$, where $A^* = (\bar{A})^t$ and $\bar{A}$ is the conjugate of the matrix $A$.
\begin{enumerate}[label=(\alph*)]
\item Show that $|\det A| \geq |a|^n$.
\item Show that if $|\det A| = |a|^n$ then $A = aI_n$.
\end{enumerate}
\end{problem}

\begin{problem}
Let integer $n \geq 2$, $a_1, a_2, \ldots, a_n$ be real numbers. The matrix
\[
A = \begin{bmatrix}
0 & 0 & 0 & \cdots & 0 & a_1 \\
1 & 0 & 0 & \cdots & 0 & a_2 \\
0 & 1 & 0 & \cdots & 0 & a_3 \\
\vdots & \vdots & \vdots & \ddots & \vdots & \vdots \\
0 & 0 & 0 & \cdots & 0 & a_{n-1} \\
0 & 0 & 0 & \cdots & 1 & a_n
\end{bmatrix}.
\]
For the linear space $V = \{X \in M_n(\mathbb{R}) \mid X^\top = X\}$, linear transformation $F : V \to V$, $F(X) = AXA^\top$. Find $\mathrm{tr}(F)$ and $\det(F)$.
\end{problem}

\begin{problem}
Let $\mathcal{S}_n^+(\mathbb{R})$ denote the set of real symmetric $n\times n$ matrices with positive eigenvalues.
Let $A \in \mathcal{S}_n^+(\mathbb{R})$. We want to show that there exists a unique matrix $B \in \mathcal{S}_n^+(\mathbb{R})$ such that
\[
B^2 = A.
\]
\begin{enumerate}[label=(\alph*)]
\item Prove existence.
\item Show that if $B \in \mathcal{S}_n^+(\mathbb{R})$ satisfies $B^2 = A$, then for every $\lambda \in \mathrm{Sp}\, A$,
\[
\mathrm{Ker}(B - \sqrt{\lambda} I_n) \subset \mathrm{Ker}(A - \lambda I_n)
\]
and then
\[
\mathrm{Ker}(B - \sqrt{\lambda} I_n) = \mathrm{Ker}(A - \lambda I_n).
\]
\item Conclude uniqueness.
\end{enumerate}
\end{problem}

\begin{problem}
Let $\mathcal{S}_n^+(\mathbb{R})$ denote the set of real symmetric $n\times n$ matrices with positive eigenvalues. Let $S \in \mathcal{S}_n^+(\mathbb{R})$.
\begin{enumerate}[label=(\alph*)]
\item Show that there exists a matrix $A \in \mathcal{S}_n^+(\mathbb{R})$ that is a polynomial in $S$ satisfying
\[
A^2 = S.
\]
\item Let $B \in \mathcal{S}_n^+(\mathbb{R})$ satisfy $B^2 = S$. Show that $B$ commutes with $A$, and hence $B = A$.
\end{enumerate}
\end{problem}

\begin{problem}
Let
\[
A = \begin{pmatrix}
10 & -9 & 0 & 0 & 0 & 0 & 0 & -9 \\
-9 & 10 & -9 & 0 & 0 & 0 & 0 & 0 \\
0 & -9 & 10 & -9 & 0 & 0 & 0 & 0 \\
0 & 0 & -9 & 10 & -9 & 0 & 0 & 0 \\
0 & 0 & 0 & -9 & 10 & -9 & 0 & 0 \\
0 & 0 & 0 & 0 & -9 & 10 & -9 & 0 \\
0 & 0 & 0 & 0 & 0 & -9 & 10 & -9 \\
-9 & 0 & 0 & 0 & 0 & 0 & -9 & 10
\end{pmatrix}.
\]
Let the eigenvalues of $A$ be $\lambda_1 \leq \cdots \leq \lambda_8$, find $\lambda_6$.
\end{problem}

\begin{problem}
Let $A$ be a real symmetric positive matrix of size $n$. For $\alpha > 0$, let
\[
\mathcal{S}_\alpha = \{M \in \mathcal{S}_n(\mathbb{R}) \mid \mathrm{Sp}\, M \subset \mathbb{R}_+ \text{ and } \det(M) \geq \alpha\}.
\]
The goal is to prove the formula:
\[
\inf_{M \in \mathcal{S}_\alpha} \mathrm{tr}(AM) = n(\alpha \det(A))^{1/n}.
\]
\begin{enumerate}[label=(\alph*)]
\item Prove the formula in the case $A = I_n$.
\item Show that every real symmetric positive matrix $A$ can be written as $A = {}^tPP$ for some square matrix $P$ of size $n$.
\item Prove the formula.
\item Is the result still true if $\alpha = 0$?
\end{enumerate}
\end{problem}

\begin{problem}
Let $m, n, k$ be positive integers. Define a linear transformation $f : M_{m\times n} \to M_{k\times n}$. Let
\[
W = \{L \mid L(AB) = L(A)B, \; \forall A \in M_{m\times n}, B \in M_{n\times n}\}.
\]
Find $\dim W$.
\end{problem}

\begin{problem}
Let $n$ be a positive integer and $C$ be the $n\times n$ matrix with entries
\[
C_{k\ell} = \cos\left(\frac{(k+\ell)\pi}{n}\right) \quad \text{for } 0 \leq k, \ell \leq n-1.
\]
Calculate $\det C$.
\end{problem}

\begin{problem}
Let $A \in \mathcal{M}_n(\mathbb{C})$ such that $A^*A^2 = AA^*$. Prove that $A^2 = A$. (Here we denote by $A^*$ the conjugate transpose of $A$.)
\end{problem}

\begin{problem}
Determine the matrices $A \in M_n(\mathbb{C})$, $n\geq 2$, with the property that $\mathrm{rank}(A^2) + \mathrm{rank}(B^2) \geq 2\,\mathrm{rank}(AB)$ for every $B \in M_n(\mathbb{C})$.
\end{problem}

\begin{problem}
Let $A$ be an $n\times n$ matrix such that $A^4 = I_n$. Prove that $A^2 + (A+I_n)^2$ and $A^2 + (A-I_n)^2$ are invertible.
\end{problem}

\begin{problem}
Let $p$ be an odd prime number and $A \in \mathcal{M}_p(\mathbb{Q})$ a matrix such that $\det(A^p + I_p) = 0$ and $\det(A + I_p) \neq 0$. Prove that:
\begin{enumerate}[label=(\alph*)]
\item $\mathrm{Tr}(A)$ is an eigenvalue of $A + I_p$.
\item $\det(A+I_p) - \det(A - I_p) = (p-1)\mathrm{Tr}(A) + 2$.
\end{enumerate}
\end{problem}

\begin{problem}
Let $k, n$ be natural numbers, $x_1, x_2, \ldots, x_k$ be distinct complex numbers and the matrix $A \in \mathcal{M}_n(\mathbb{C})$ such that $(A - x_1 I_n)(A - x_2 I_n) \cdots (A - x_k I_n) = O_n$. Prove that
\[
\mathrm{rank}(A - x_1 I_n) + \mathrm{rank}(A - x_2 I_n) + \cdots + \mathrm{rank}(A - x_k I_n) = n(k-1).
\]
\end{problem}

\begin{problem}
Let $K$ be a field and let $n \geq 1$. Let $A, B \in \mathcal{M}_n(K)$ such that $[A,B]$ commutes with $A$ or $B$.

If $\mathrm{char}\, K = 0$ or $\mathrm{char}\, K > n$ then it is known that $[A,B]$ is nilpotent, i.e., $[A,B]^n = 0$.

Prove that this result no longer holds if $0 < \mathrm{char}\, K \leq n$. (Here by $[\cdot,\cdot]$ we mean the commutator, $[A,B] = AB - BA$.)
\end{problem}

\begin{problem}
Prove that, for a matrix $A \in \mathcal{M}_n(\mathbb{C})$, the next assertions are equivalent:
\begin{enumerate}[label=(\alph*)]
\item $A^2 = A$;
\item $\mathrm{rank}\,A + \mathrm{rank}(I_n - A) = n$.
\end{enumerate}
\end{problem}

\begin{problem}
Consider matrices $A, B \in \mathcal{M}_n(\mathbb{R})$.
\begin{enumerate}[label=(\alph*)]
\item Show that there exists $a > 0$ such that, for every $\varepsilon \in (-a,a)$, $\varepsilon \neq 0$, the matrix equation
\[
AX + \varepsilon X = B, \qquad X \in \mathcal{M}_n(\mathbb{R}),
\]
has a unique solution $X(\varepsilon) \in \mathcal{M}_n(\mathbb{R})$.
\item If $B^2 = I_n$ and $A$ is diagonalizable, show that
\[
\lim_{\varepsilon \to 0} \varepsilon\, \mathrm{Tr}(BX(\varepsilon)) = n - \mathrm{rank}\, A,
\]
where $\mathrm{Tr}$ denotes the trace of a matrix.
\end{enumerate}
\end{problem}

\begin{problem}
\begin{enumerate}[label=(\alph*)]
\item Show that for every matrix $X \in \mathcal{M}_2(\mathbb{C})$ there exists a matrix $Y \in \mathcal{M}_2(\mathbb{C})$ such that $Y^3 = X^2$.
\item Show that there exists a matrix $X \in \mathcal{M}_3(\mathbb{C})$ such that for every matrix $Y \in \mathcal{M}_3(\mathbb{C})$, $Y^3 \neq X^2$.
\end{enumerate}
\end{problem}

\begin{problem}
Let $n, k \in \mathbb{N}^*$ and let $A_1, \ldots, A_k \in \mathcal{M}_n(\mathbb{R})$ be idempotent matrices. Show that
\[
\sum_{i=1}^{k} (n - \mathrm{rank}(A_i)) \geq \mathrm{rank}\left(I_n - \prod_{i=1}^{k} A_i\right).
\]
\end{problem}

\begin{problem}
Let $H_k = \sum_{j=1}^{k} 1/j$, and let $D_n$ be the determinant of the $(n+1)\times(n+1)$ Hankel matrix with $(i,j)$ entry $H_{i+j+1}$ for $1 \leq i,j \leq n$. (Thus, $D_1 = -5/12$ and $D_2 = 1/216$.) Show that for $n \geq 1$,
\[
D_n = \frac{\prod_{i=1}^{n} i!^4}{\prod_{i=1}^{2n+1} i!} \cdot \sum_{j=0}^{n} \frac{(-1)^j (n+j+1)! H_{j+1}}{j!(j+1)!(n-j)!}.
\]
\textit{Hint.} Let $H$ be the $(n+1)\times(n+1)$ Hilbert matrix with $(i,j)$ entry $1/(i+j+1)$ for $0 \leq i,j \leq n$. Show that
\[
\det(H) = \frac{\prod_{i=1}^{n} i!^4}{\prod_{i=1}^{2n+1} i!}.
\]
\end{problem}

\begin{problem}
Let $A$ be a positive-definite $n \times n$ Hermitian matrix with minimum eigenvalue $\lambda$ and maximum eigenvalue $\Lambda$. Show that
\[
\left( \frac{n}{\mathrm{tr}((A+\lambda I)^{-1})} - \lambda \right)^n \leq \det(A) \leq \left( \frac{n}{\mathrm{tr}((A+\Lambda I)^{-1})} - \Lambda \right)^n.
\]
\textit{Hint:} Notice that $\mathrm{tr}((A+\lambda I)^{-1}) = \sum_{k=1}^{n} 1/(\lambda + \lambda_k)$, where $\lambda_1, \ldots, \lambda_n$ denote the eigenvalues of $A$.
\end{problem}

\begin{problem}
Let $A$ be an $r \times r$ matrix with distinct eigenvalues $\lambda_1, \ldots, \lambda_n$. For $n \geq 0$, let $a(n)$ be the trace of $A^n$. Let $H(n)$ be the $r \times r$ Hankel matrix with $(i,j)$ entry $a(i+j+n-2)$. Show that
\[
\lim_{n \to \infty} |\det H(n)|^{1/n} = \prod_{k=1}^{r} |\lambda_k|.
\]
\end{problem}

\begin{problem}
Let $M$ be the set of all $3 \times 3$ matrices with complex entries. Suppose $A, B \in M$ with $AB = BA$ and $\mathrm{Tr}(A^k) = \mathrm{Tr}(B^k)$ for $k \in \{1,2,3\}$. Suppose further that $n \geq 4$ and that $A^n - B^n$ is invertible. Prove that there exist complex numbers $\alpha, \beta, \gamma$ such that
\[
A^{2n} + B^{2n} + A^n B^n + \alpha A^n + \beta B^n + \gamma I = 0.
\]
\end{problem}

\begin{problem}
Let $n$ be a positive integer, and let $A = (a_{ij})_{i,j=1}^n$ be the $n\times n$ matrix with $a_{ii} = 2$, $a_{i,i\pm 1} = -1$, and $a_{ij} = 0$ otherwise; that is,
\[
A = \begin{pmatrix}
2 & -1 & 0 & 0 & \cdots & 0 & 0 & 0 \\
-1 & 2 & -1 & 0 & \cdots & 0 & 0 & 0 \\
0 & -1 & 2 & -1 & \cdots & 0 & 0 & 0 \\
0 & 0 & -1 & 2 & \cdots & 0 & 0 & 0 \\
\vdots & \vdots & \vdots & \vdots & \ddots & \vdots & \vdots & \vdots \\
0 & 0 & 0 & 0 & \cdots & 2 & -1 & 0 \\
0 & 0 & 0 & 0 & \cdots & -1 & 2 & -1 \\
0 & 0 & 0 & 0 & \cdots & 0 & -1 & 2
\end{pmatrix}.
\]
Prove that every eigenvalue of $A$ is a positive real number.
\end{problem}

\begin{problem}
Let $A$ be a real symmetric $n \times n$ matrix with nonnegative entries. Prove that $A$ has an eigenvector with nonnegative entries.
\end{problem}

\begin{problem}
Let $A_n$ be the $n \times n$ matrix whose entries $a_{jk}$ are given by
\[
a_{jk} = \begin{cases} 1 & \text{if } |j-k| = 1 \\ 0 & \text{otherwise.} \end{cases}
\]
Prove that the eigenvalues of $A$ are symmetric with respect to the origin.
\end{problem}

\begin{problem}
Let $A = (a_{ij})_{i,j=1}^n$ be a real $n \times n$ matrix with nonnegative entries such that
\[
\sum_{j=1}^{n} a_{ij} = 1 \qquad (1 \leq i \leq n).
\]
Prove that no eigenvalue of $A$ has absolute value greater than $1$.
\end{problem}

\begin{problem}
Let $S$ be a special orthogonal $n \times n$ matrix, a real $n \times n$ matrix satisfying $S^t S = I$ and $\det(S) = 1$.
\begin{enumerate}[label=\arabic*.]
\item Prove that if $n$ is odd then $1$ is an eigenvalue of $S$.
\item Prove that if $n$ is even then $1$ need not be an eigenvalue of $S$.
\end{enumerate}
\end{problem}

\begin{problem}
Let $A$ and $B$ be two $n \times n$ self-adjoint (i.e., Hermitian) matrices over $\mathbb{C}$ such that all eigenvalues of $A$ lie in $[a,a']$ and all eigenvalues of $B$ lie in $[b,b']$. Show that all eigenvalues of $A+B$ lie in $[a+b, a'+b']$.
\end{problem}

\begin{problem}
Let $k$ be real, $n$ an integer $\geq 2$, and let $A = (a_{ij})$ be the $n \times n$ matrix such that all diagonal entries $a_{ii} = k$, all entries $a_{i,i\pm 1}$ immediately above or below the diagonal equal $1$, and all other entries equal $0$. For example, if $n=5$,
\[
A = \begin{pmatrix}
k & 1 & 0 & 0 & 0 \\
1 & k & 1 & 0 & 0 \\
0 & 1 & k & 1 & 0 \\
0 & 0 & 1 & k & 1 \\
0 & 0 & 0 & 1 & k
\end{pmatrix}.
\]
Let $\lambda_{\min}$ and $\lambda_{\max}$ denote the smallest and largest eigenvalues of $A$, respectively. Show that $\lambda_{\min} \leq k-1$ and $\lambda_{\max} \geq k+1$.
\end{problem}

\begin{problem}
If $A$ is an $n \times n$ matrix such that $A^3 = pA^2 + qA + rI$, determine $(A-I)^{-1}$ assuming it exists.
\end{problem}

\begin{problem}
Let $r,s,t$ be non-zero linear transformations on a finite-dimensional vector space $V$ such that $r\circ t\circ r=0$. Let $p=r\circ s$ and $q=r\circ(s+t)$, and suppose that the minimum polynomials of $p,q$ are $p(X),q(X)$ respectively. Prove that (with composites written as products)
\begin{enumerate}[label=(\arabic*)]
\item $q^n=p^{n-1}q$ and $p^n=q^{n-1}p$ for all $n\ge 1$;
\item $p(X)$ and $q(X)$ are divisible by $X$;
\item $q$ satisfies $Xp(X)=0$, and $p$ satisfies $Xq(X)=0$.
\end{enumerate}
Deduce that one of the following holds:
\begin{enumerate}[label=(\roman*)]
\item $p(X)=q(X)$;
\item $p(X)=Xq(X)$;
\item $q(X)=Xp(X)$.
\end{enumerate}
\end{problem}

\begin{problem}
A $3\times 3$ complex matrix $M$ is said to be \emph{magic} if every row sum, every column sum, and both diagonal sums are equal to some $\vartheta\in\mathbb{C}$.

If $M$ is magic, prove that $\vartheta=3m_{22}$. Deduce that, given $\alpha,\beta,\gamma\in\mathbb{C}$, there is a unique magic matrix $M(\alpha,\beta,\gamma)$ such that
\[
m_{22}=\alpha,\qquad m_{11}=\alpha+\beta,\qquad m_{31}=\alpha+\gamma.
\]
Show that $\{M(\alpha,\beta,\gamma)\mid \alpha,\beta,\gamma\in\mathbb{C}\}$ is a subspace of $\mathrm{Mat}_{3\times 3}(\mathbb{C})$ and that
\[
B=\{M(1,0,0),\ M(0,1,0),\ M(0,0,1)\}
\]
is a basis of this subspace.

If $f:\mathbb{C}^3\to\mathbb{C}^3$ represents $M(\alpha,\beta,\gamma)$ relative to the canonical basis $\{e_1,e_2,e_3\}$, show that $e_1+e_2+e_3$ is an eigenvector of $f$. Determine the matrix of $f$ relative to the basis $\{e_1+e_2+e_3,e_2,e_3\}$. Hence find the eigenvalues of $M(\alpha,\beta,\gamma)$.
\end{problem}

\end{document}
